% =====================================================
% CHAPITRE : INTERFACE FRONTEND
% À inclure dans rapport_pfe.tex via % =====================================================
% CHAPITRE : INTERFACE FRONTEND
% À inclure dans rapport_pfe.tex via % =====================================================
% CHAPITRE : INTERFACE FRONTEND
% À inclure dans rapport_pfe.tex via % =====================================================
% CHAPITRE : INTERFACE FRONTEND
% À inclure dans rapport_pfe.tex via \input{chapitre_frontend}
% =====================================================

\chapter{Interface Frontend de la Plateforme}

\section*{Introduction du chapitre}
\addcontentsline{toc}{section}{Introduction du chapitre}

La conception de l'interface frontend a répondu à un défi central~: mettre entre les mains d'auditeurs, dont la majorité ne sont pas informaticiens, un outil d'intelligence artificielle puissant, sans que la complexité technique ne transparaisse dans leur expérience quotidienne. L'objectif était donc de construire une interface \textbf{professionnelle, intuitive et rassurante}, fidèle à l'image de marque de PwC, tout en rendant lisibles les résultats d'algorithmes d'apprentissage automatique sophistiqués.

Ce chapitre présente les choix technologiques retenus, l'architecture du code, les interfaces utilisateur développées, la manière dont le frontend communique avec les deux backends distincts, ainsi que les mécanismes d'authentification, de contrôle d'accès et de déploiement.

% =====================================================
\section{Stack Technologique}
% =====================================================

\subsection{Framework principal~: Next.js 15 avec App Router}

Le choix de \textbf{Next.js 15} (version la plus récente au moment du développement) avec le système \textbf{App Router} constitue la colonne vertébrale de toute l'interface. Ce framework, construit au-dessus de React 19, offre plusieurs avantages déterminants pour ce type d'application d'entreprise~:

\begin{itemize}[leftmargin=1.5em]
  \item Le \textbf{rendu hybride} permet de combiner rendu côté serveur (pour les données critiques comme les missions et les résultats d'analyse) et rendu côté client (pour les interactions dynamiques comme les filtres et les graphiques), améliorant à la fois la performance perçue et la sécurité.
  \item L'\textbf{App Router} restructure l'organisation du code autour de la notion de layout partagé, de chargement progressif et de routes groupées, ce qui se traduit par une navigation fluide sans rechargement de page.
  \item La gestion intégrée des \textbf{rewrites de proxy} redirige toutes les requêtes vers le backend FastAPI sans exposer l'URL réelle de ce dernier au navigateur, évitant les problèmes CORS et renforçant la sécurité de l'architecture.
\end{itemize}

\textbf{TypeScript 5} est utilisé en mode strict sur l'ensemble du projet. Chaque donnée échangée entre le frontend et les backends est typée explicitement --- missions, utilisateurs, résultats d'analyse, anomalies --- ce qui élimine une classe entière de bugs liés aux données inattendues ou manquantes lors du développement.

\subsection{Stylisation~: TailwindCSS et shadcn/ui}

\textbf{TailwindCSS 3.4} est le système de mise en page et de style. Il adopte une approche utilitaire~: les styles sont composés directement dans le JSX via des classes courtes (\texttt{flex}, \texttt{gap-4}, \texttt{text-lg}, etc.), ce qui accélère le développement tout en garantissant la cohérence visuelle. La palette de couleurs est entièrement configurée pour respecter l'identité PwC~:

\begin{itemize}[leftmargin=1.5em]
  \item \textbf{Orange primaire \texttt{\#D04A02}}~: boutons d'action principaux, accents visuels
  \item \textbf{Bleu foncé \texttt{\#293854}}~: en-têtes et éléments de navigation
  \item \textbf{Niveaux de risque}~: Rouge \texttt{\#C00000} (CRITIQUE), Orange \texttt{\#D04A02} (ÉLEVÉ), Vert \texttt{\#008246} (FAIBLE)
\end{itemize}

\textbf{shadcn/ui} est la bibliothèque de composants retenue. Contrairement à des solutions classiques comme Material UI, shadcn/ui fournit des composants \textit{headless} basés sur Radix UI~: sans style imposé mais avec une accessibilité garantie (navigation clavier, attributs ARIA, compatibilité lecteurs d'écran). Les composants intégrés comprennent Dialog (modales), Tabs (onglets de résultats), Progress (barre de progression d'upload), Select, Dropdown, Badge, Card, Button et Skeleton (placeholders de chargement).

\subsection{Gestion des données~: TanStack React Query 5}

\textbf{TanStack React Query} gère l'ensemble de la récupération, mise en cache et synchronisation des données provenant des APIs. Son fonctionnement est fondamental pour la fluidité de l'expérience utilisateur~:

\begin{itemize}[leftmargin=1.5em]
  \item Les données (missions, utilisateurs, analyses) sont mises en cache \textbf{2 minutes}~: naviguer entre pages ne déclenche pas d'appels réseau répétés.
  \item Après une modification (création de mission, upload de fichier), le cache est \textbf{invalidé et rechargé automatiquement en arrière-plan}, sans que l'utilisateur ait à rafraîchir la page.
  \item Les mutations gèrent automatiquement les états de chargement et d'erreur, déclenchant des notifications via la bibliothèque Sonner.
\end{itemize}

\subsection{Bibliothèques complémentaires}

Le tableau \ref{tab:frontend-libs} récapitule les bibliothèques complémentaires intégrées au projet.

\begin{table}[H]
\centering
\caption{Bibliothèques frontend complémentaires}
\label{tab:frontend-libs}
\begin{tabular}{>{\bfseries}p{4.5cm} p{9.5cm}}
\toprule
Bibliothèque & Usage \\
\midrule
Recharts 2.13 & Graphiques interactifs (camembert, histogramme, barres) basés sur D3.js avec une API React native \\
Lucide React 0.468 & Bibliothèque d'icônes SVG (tree-shaken, seules les icônes utilisées sont embarquées dans le bundle) \\
Sonner 1.7 & Notifications toast positionnées en haut à droite (succès, erreur, information, 4 secondes) \\
React Hook Form 7.54 & Gestion d'état des formulaires sans re-rendu inutile \\
Zod 3.23 & Validation de schéma à l'exécution (email \texttt{@pwc.com}, dates, formats) \\
date-fns 4.1 & Formatage des dates en français (\texttt{dd/MM/yyyy HH:mm}) \\
Axios 1.7 & Client HTTP avec intercepteurs pour l'injection automatique du token JWT \\
Prisma ORM 5.22 & Client base de données avec typage auto-généré pour PostgreSQL \\
bcryptjs 3.0 & Hachage des mots de passe (12 rounds, résistant aux attaques par force brute) \\
jose 5.9 & Génération et vérification des tokens JWT (algorithme HS256) \\
uuid 11.0 & Génération d'identifiants uniques côté client \\
\bottomrule
\end{tabular}
\end{table}

% =====================================================
\section{Architecture du Code}
% =====================================================

L'organisation du code suit les conventions Next.js App Router, structurées en couches logiques clairement séparées~:

\begin{lstlisting}[caption={Structure des dossiers du frontend}, label={lst:frontend-structure}]
frontend/src/
  app/                        # Routes et pages (App Router)
    (auth)/login/             # Page de connexion publique
    (dashboard)/              # Toutes les pages authentifiees
      dashboard/              # Tableau de bord principal
      missions/               # Liste des missions
        [id]/                 # Detail d'une mission
          analysis/           # Assistant d'analyse (wizard 4 etapes)
      history/                # Historique des analyses
      reports/                # Rapports PDF/DOCX generes
      audit-trail/            # Piste d'audit complete
      admin/users/            # Gestion des utilisateurs
    api/                      # Routes API internes Next.js
      auth/                   # Login, logout, me
      missions/               # CRUD missions
      users/                  # Gestion utilisateurs
      analysis-runs/          # Persistance des analyses
      audit-logs/             # Journal d'audit
  components/                 # Composants reutilisables
    analysis/                 # Wizard, tableau de resultats, KPI
    charts/                   # Graphiques Recharts
    datasets/                 # Upload et gestion de fichiers
    missions/                 # Cartes et formulaires de mission
    reports/                  # Generation PDF/DOCX
    explanations/             # Affichage SHAP, LIME, LLM
    layout/                   # Navbar, TopBar
    ui/                       # Primitives shadcn/ui
  lib/
    api/                      # Services d'appel aux APIs
    hooks/                    # Hooks personnalises
    store/                    # Stores locaux (mode demo)
  providers/                  # Providers React (Query, Theme, Language)
  types/                      # Interfaces TypeScript partages
  messages/                   # Traductions (fr.json, en.json)
\end{lstlisting}

% =====================================================
\section{Routes et Interfaces Utilisateur}
% =====================================================

\subsection{Page de Connexion (\texttt{/login})}

La page d'entrée est volontairement sobre et professionnelle. Elle présente le logo PwC et un formulaire à deux champs (email et mot de passe). La validation s'effectue à deux niveaux~:

\begin{itemize}[leftmargin=1.5em]
  \item \textbf{Côté client (Zod)}~: le format email est vérifié et le domaine \texttt{@pwc.com} est obligatoire. Un message d'erreur s'affiche immédiatement sans solliciter le serveur.
  \item \textbf{Côté serveur}~: le mot de passe est comparé au hash bcrypt stocké en base de données (12 rounds, environ 100\,ms par vérification --- délibérément lent pour résister aux attaques par force brute).
\end{itemize}

\subsection{Tableau de Bord Principal (\texttt{/dashboard})}

Le tableau de bord est la première page visible après la connexion. Il offre une vue synthétique de l'activité en cours~: cartes KPI (missions actives, analyses lancées, anomalies détectées), activité récente, et accès rapide aux missions. Le contenu s'adapte automatiquement au \textbf{rôle de l'utilisateur connecté}~: un auditeur ne voit que ses missions assignées, tandis qu'un manager dispose d'une vue d'ensemble complète.

\subsection{Gestion des Missions (\texttt{/missions})}

La liste des missions est présentée sous forme de cartes, chacune affichant le nom de la mission, l'entreprise cliente, le type (commissariat aux comptes, audit interne, due diligence, etc.), les dates, le statut (en préparation, en cours, terminée, archivée) et les auditeurs assignés. Une barre de recherche et des filtres combinés par statut et type permettent de naviguer rapidement.

Le bouton \textit{Nouvelle mission} --- visible uniquement pour les managers, partenaires et administrateurs --- ouvre une modale de création avec un formulaire complet incluant l'affectation d'auditeurs depuis la liste des utilisateurs actifs.

\subsection{Détail d'une Mission (\texttt{/missions/[id]})}

La page de détail consolide~:

\begin{itemize}[leftmargin=1.5em]
  \item Les \textbf{métadonnées de la mission}~: informations générales, équipe, calendrier.
  \item La \textbf{section Datasets}~: liste des fichiers CSV uploadés avec drag-and-drop (\texttt{UploadDropzone}), barre de progression, possibilité de remplacer ou supprimer un fichier (jusqu'à 500\,Mo).
  \item Le \textbf{bouton Lancer l'analyse}~: redirection vers l'assistant multi-étapes.
\end{itemize}

\subsection{Assistant d'Analyse (\texttt{/missions/[id]/analysis})}

C'est l'interface la plus sophistiquée de la plateforme. Elle guide l'utilisateur à travers quatre étapes structurées~:

\begin{description}[leftmargin=1.5em, style=nextline]
  \item[Étape 1 --- Upload du fichier CSV]
    L'utilisateur dépose ou sélectionne le fichier de transactions. Un indicateur de santé du backend ML s'affiche (modèles chargés, LLM disponible) pour rassurer l'utilisateur avant de démarrer.

  \item[Étape 2 --- Sélection du modèle]
    Trois configurations sont proposées avec leur description~:
    \begin{itemize}[leftmargin=1em]
      \item \textbf{Combiné (recommandé)}~: les trois modèles votent en ensemble pour une décision plus robuste.
      \item \textbf{XGBoost uniquement}~: rapide et précis sur les patterns de fraude connus.
      \item \textbf{AutoEncoder uniquement}~: détecte des anomalies inédites, sans supervision.
    \end{itemize}

  \item[Étape 3 --- Exécution de l'analyse]
    Une carte récapitulative s'affiche avant le lancement. Pendant l'analyse, une barre de progression et des messages d'état informent l'utilisateur de chaque phase~: détection du schéma, ingénierie des features, inférence. Le timeout est fixé à 120\,secondes.

  \item[Étape 4 --- Tableau de bord des résultats]
    Les résultats sont organisés en quatre onglets~:
    \begin{itemize}[leftmargin=1em]
      \item \textbf{Graphiques}~: camembert normal/anomalie, histogramme des scores avec ligne de seuil, barres par catégorie.
      \item \textbf{Transactions}~: tableau interactif et filtrable de toutes les transactions avec score de risque et niveau (CRITIQUE/ÉLEVÉ/FAIBLE).
      \item \textbf{Explication}~: pour une transaction sélectionnée, résumé LLM en français, valeurs SHAP, règles LIME, recommandations d'actions.
      \item \textbf{Rapport}~: génération et téléchargement du rapport en PDF ou DOCX.
    \end{itemize}
\end{description}

\subsection{Autres Pages}

\begin{itemize}[leftmargin=1.5em]
  \item \textbf{Historique (\texttt{/history})}~: liste de toutes les analyses précédentes avec filtres, permettant de retrouver et revisualiser des résultats sans relancer une analyse.
  \item \textbf{Rapports (\texttt{/reports})}~: centralise tous les rapports générés, avec téléchargement et archivage.
  \item \textbf{Piste d'audit (\texttt{/audit-trail})}~: journal complet de toutes les actions (connexions, uploads, analyses, rapports) avec horodatage et identité de l'acteur. Réservé aux managers et administrateurs.
  \item \textbf{Gestion des utilisateurs (\texttt{/admin/users})}~: création, modification et désactivation de comptes, assignation de rôles. Réservé aux administrateurs.
\end{itemize}

% =====================================================
\section{Architecture API et Communication avec le Backend}
% =====================================================

\subsection{Double couche d'API}

Le frontend communique avec deux backends distincts selon une architecture en deux niveaux~:

\begin{description}[leftmargin=1.5em, style=nextline]
  \item[Couche 1 --- Next.js API Routes (PostgreSQL)]
    Un backend léger intégré à Next.js, accessible via \texttt{/api/*}, gère toute la persistance métier~: utilisateurs, missions, datasets, historique d'analyses, piste d'audit. Ces routes s'appuient sur Prisma ORM pour interagir avec une base PostgreSQL.

  \item[Couche 2 --- FastAPI ML Backend]
    Le service Python qui héberge les modèles de machine learning, accessible via \texttt{/ml/api/*}. Le frontend ne communique jamais directement avec ce service~: les requêtes transitent par un \textbf{proxy configuré dans \texttt{next.config.ts}}, qui redirige les appels \texttt{/ml/*} vers \texttt{http://localhost:8000}. Ce mécanisme élimine les problèmes CORS et masque l'URL réelle du backend ML au navigateur.
\end{description}

\subsection{Services API}

Le tableau \ref{tab:api-ml} présente les endpoints ML consommés par le frontend, et le tableau \ref{tab:api-app} les endpoints internes.

\begin{table}[H]
\centering
\caption{Endpoints FastAPI ML consommés par le frontend}
\label{tab:api-ml}
\begin{tabular}{p{3.5cm} p{3cm} p{7cm}}
\toprule
\textbf{Service} & \textbf{Méthode / Route} & \textbf{Usage} \\
\midrule
\texttt{predict} & POST \texttt{/api/predict} & Upload CSV, pipeline complet (détection schéma, feature engineering, inférence) \\
\texttt{explain} & GET \texttt{/api/explain/\{tx\_id\}} & SHAP + LIME + explication LLM pour une transaction \\
\texttt{explainBatch} & POST \texttt{/api/explain/batch} & Explications en lot (LIME désactivé pour performance) \\
\texttt{getModels} & GET \texttt{/api/models} & Métriques des modèles (recall, précision, F1, ROC-AUC) \\
\texttt{profile} & POST \texttt{/api/profile} & Profiling qualité du dataset sans prédiction \\
\texttt{health} & GET \texttt{/api/health} & Santé du backend (modèles chargés, LLM disponible) \\
\texttt{generatePDF} & POST \texttt{/api/report} & Génération du rapport PDF \\
\texttt{generateDOCX} & POST \texttt{/api/report/docx} & Génération du rapport Word \\
\bottomrule
\end{tabular}
\end{table}

\begin{table}[H]
\centering
\caption{Endpoints Next.js internes (PostgreSQL)}
\label{tab:api-app}
\begin{tabular}{p{4cm} p{3.5cm} p{6cm}}
\toprule
\textbf{Service} & \textbf{Route} & \textbf{Usage} \\
\midrule
Auth & \texttt{/api/auth/login} & Émission du JWT \\
& \texttt{/api/auth/logout} & Invalidation de session \\
& \texttt{/api/auth/me} & Utilisateur courant \\
Missions & \texttt{/api/missions} & Liste / création \\
& \texttt{/api/missions/[id]} & Lecture / modification / suppression \\
Datasets & \texttt{/api/missions/[id]/datasets} & Liste / upload / remplacement / suppression \\
Utilisateurs & \texttt{/api/users} & Liste et détail \\
Analyses & \texttt{/api/analysis-runs} & Persistance de l'historique \\
Audit & \texttt{/api/audit-logs} & Enregistrement des actions \\
\bottomrule
\end{tabular}
\end{table}

\subsection{Intercepteurs Axios}

Les deux clients HTTP sont configurés avec des intercepteurs~:

\begin{itemize}[leftmargin=1.5em]
  \item \textbf{Intercepteur de requête}~: injecte automatiquement l'en-tête \texttt{Authorization: Bearer \{token\}} à chaque appel, en lisant le token depuis le \texttt{localStorage}.
  \item \textbf{Intercepteur de réponse}~:
    \begin{itemize}[leftmargin=1em]
      \item Réponse \textbf{401} (token expiré ou invalide) → nettoyage du \texttt{localStorage} et redirection vers \texttt{/login}.
      \item \textbf{Erreur réseau} → une tentative de relance automatique.
      \item \textbf{Erreur 4xx} (erreur client) → pas de relance.
    \end{itemize}
\end{itemize}

% =====================================================
\section{Authentification et Contrôle d'Accès}
% =====================================================

\subsection{Flux JWT}

L'authentification repose sur des \textbf{tokens JWT} (JSON Web Tokens) signés avec l'algorithme HS256. Le flux est le suivant~:

\begin{enumerate}[leftmargin=1.5em]
  \item L'utilisateur soumet son email (\texttt{@pwc.com} obligatoire) et son mot de passe.
  \item Le serveur vérifie l'email et compare le mot de passe au hash bcrypt (12 rounds).
  \item Un JWT est généré avec les informations de l'utilisateur (id, email, nom, rôle) et une durée de validité de \textbf{8 heures}.
  \item Le token est retourné au client et stocké dans le \texttt{localStorage} (\texttt{pwc\_token}).
  \item À chaque requête ultérieure, l'intercepteur Axios injecte le token dans l'en-tête \texttt{Authorization}.
  \item Si le token expire, le serveur répond 401 → déconnexion automatique et redirection vers \texttt{/login}.
\end{enumerate}

\subsection{Rôles et Permissions}

Quatre rôles définissent les droits d'accès dans la plateforme~:

\begin{description}[leftmargin=1.5em, style=nextline]
  \item[Auditeur] Accède uniquement aux missions qui lui sont assignées. Peut uploader des fichiers CSV, lancer des analyses et commenter des anomalies.
  \item[Manager] Peut créer et modifier des missions, assigner des auditeurs, approuver des rapports. Accès à toutes les missions.
  \item[Partenaire] Peut créer des missions et consulter les rapports générés.
  \item[Administrateur] Accès complet à toutes les fonctionnalités, y compris la création et la modification des comptes utilisateurs.
\end{description}

Le hook personnalisé \texttt{usePermissions()} est utilisé dans chaque composant pour afficher ou masquer conditionnellement les éléments d'interface selon le rôle de l'utilisateur connecté.

% =====================================================
\section{Déploiement des Modèles côté Frontend}
% =====================================================

\subsection{Sélection des modèles}

L'étape 2 de l'assistant d'analyse (\texttt{AnalysisWizard.tsx}) présente trois configurations de modèles. Le choix de l'utilisateur est transmis en paramètre (\texttt{model\_type}) dans le corps de la requête à \texttt{POST /ml/api/predict}, avec les valeurs possibles \texttt{combined}, \texttt{xgboost} ou \texttt{autoencoder}.

\subsection{Orchestration de l'analyse}

Quand l'utilisateur déclenche une analyse~:

\begin{enumerate}[leftmargin=1.5em]
  \item Le fichier CSV est envoyé en \textit{multipart/form-data} avec le type de modèle choisi.
  \item Une barre de progression s'affiche --- le timeout est fixé à 120 secondes pour accommoder les grands fichiers.
  \item Le backend FastAPI exécute séquentiellement~: détection du schéma, profiling qualité, ingénierie des 14 features dérivées, puis inférence des modèles sélectionnés.
  \item La réponse contient l'ensemble des transactions avec leurs scores et niveaux de risque.
  \item Si la détection de schéma est incomplète (confiance sous le seuil), un \textbf{avertissement} est affiché pour informer l'auditeur de l'impact possible sur la précision.
\end{enumerate}

\subsection{Explications à la demande}

Les explications (SHAP, LIME, résumé LLM) ne sont pas calculées en masse au moment de l'analyse~: elles sont générées \textbf{à la demande}, pour la transaction que l'auditeur sélectionne dans le tableau. Cela évite une surcharge inutile du backend pour des transactions qui ne seront jamais inspectées en détail. Pour la génération de rapport, un appel batch (\texttt{POST /ml/api/explain/batch}) permet de récupérer les explications des anomalies principales en une seule requête, avec LIME désactivé pour des raisons de performance.

% =====================================================
\section{Internationalisation et Thème}
% =====================================================

\subsection{Interface bilingue français/anglais}

L'interface est entièrement bilingue. Les traductions sont stockées dans \texttt{src/messages/fr.json} et \texttt{en.json}. La langue est persistée dans le \texttt{localStorage} (\texttt{pwc\_locale}) et peut être changée à tout moment via un sélecteur dans la barre de navigation. Le \texttt{LanguageContext} injecte la fonction de traduction \texttt{t()} dans tous les composants qui en ont besoin, avec support de l'interpolation de variables.

\subsection{Mode sombre}

Un bouton dans la barre de navigation permet de basculer entre mode clair et mode sombre. Le \texttt{ThemeProvider} gère cela via la classe \texttt{dark} sur l'élément racine HTML. TailwindCSS applique automatiquement les variantes \texttt{dark:} définies sur chaque composant.

% =====================================================
\section{Performance et Configuration}
% =====================================================

\subsection{Optimisations de build (\texttt{next.config.ts})}

\begin{itemize}[leftmargin=1.5em]
  \item \textbf{Tree-shaking des imports baril}~: seules les icônes et composants réellement utilisés de \texttt{lucide-react}, \texttt{recharts} et \texttt{@radix-ui} sont inclus dans le bundle final.
  \item \textbf{Cache agressif des assets statiques}~: 31\,536\,000 secondes (1 an) pour les images et les fichiers Next.js compilés.
  \item \textbf{Optimisation des images}~: conversion automatique en AVIF et WebP (40 à 50\,\% plus légers que JPEG/PNG).
  \item \textbf{Compression}~: activée sur toutes les réponses HTTP.
  \item \textbf{Standalone output}~: pour Docker, le build n'embarque que les dépendances strictement nécessaires à l'exécution, produisant une image légère.
  \item \textbf{Limite de 500\,Mo}~: configurée pour les Server Actions afin de permettre l'upload de fichiers CSV volumineux.
\end{itemize}

\subsection{Schéma de base de données}

Le schéma Prisma (\texttt{prisma/schema.prisma}) définit 11 modèles interconnectés~: \texttt{User}, \texttt{Mission}, \texttt{MissionAssignment}, \texttt{Dataset}, \texttt{DatasetVersion}, \texttt{AnalysisRun}, \texttt{Anomaly}, \texttt{AnomalyComment}, \texttt{Report}, \texttt{AuditLog} et \texttt{Permission}. Des index sont définis sur les colonnes les plus fréquemment filtrées (\texttt{userId}, \texttt{missionId}, \texttt{action}, \texttt{timestamp}) pour maintenir les performances de requête à mesure que les données s'accumulent.

% =====================================================
\section{Gestion des Erreurs et Expérience Utilisateur}
% =====================================================

Chaque opération susceptible d'échouer est traitée explicitement~:

\begin{itemize}[leftmargin=1.5em]
  \item \textbf{Upload échoué}~: message descriptif formaté depuis l'erreur serveur, avec possibilité de réessayer.
  \item \textbf{Analyse échouée}~: message extrait de la réponse FastAPI, avec code d'erreur si disponible.
  \item \textbf{Token expiré}~: déconnexion propre et redirection sans page blanche ni message cryptique.
  \item \textbf{Backend indisponible}~: indicateur rouge dans l'assistant d'analyse avec message d'alerte explicite.
  \item \textbf{Schéma CSV non reconnu}~: avertissement détaillé sur les colonnes manquantes et leur impact potentiel sur la précision des modèles.
\end{itemize}

Les \textbf{notifications Sonner} (4 secondes, en haut à droite) confirment toutes les actions réussies~: mission créée, fichier uploadé, analyse terminée, rapport généré.

Les \textbf{états de chargement} sont gérés avec des squelettes (\texttt{Skeleton}) qui maintiennent la mise en page pendant les requêtes, évitant les sauts visuels qui nuisent à la perception de fiabilité de l'outil.

% =====================================================
\section*{Conclusion du chapitre}
\addcontentsline{toc}{section}{Conclusion du chapitre}
% =====================================================

Le frontend de cette plateforme est une application React de niveau professionnel, pensée pour un contexte d'entreprise exigeant. Sa cohérence repose sur une architecture claire~: Next.js 15 pour le rendu hybride et le proxy ML, TanStack Query pour la gestion de l'état serveur, shadcn/ui pour l'accessibilité, et Axios avec intercepteurs pour la sécurité des échanges.

La force de cette interface réside dans la \textbf{lisibilité du parcours utilisateur}~: du premier clic sur \textit{Connexion} jusqu'au téléchargement du rapport PDF, chaque étape est guidée et protégée par des validations, des messages d'erreur clairs et des indicateurs de progression. La complexité des modèles d'intelligence artificielle est totalement abstraite~: l'auditeur choisit un fichier, sélectionne une configuration, lance l'analyse, et reçoit des résultats exploitables en quelques dizaines de secondes, accompagnés d'explications rédigées en français naturel.

% =====================================================

\chapter{Interface Frontend de la Plateforme}

\section*{Introduction du chapitre}
\addcontentsline{toc}{section}{Introduction du chapitre}

La conception de l'interface frontend a répondu à un défi central~: mettre entre les mains d'auditeurs, dont la majorité ne sont pas informaticiens, un outil d'intelligence artificielle puissant, sans que la complexité technique ne transparaisse dans leur expérience quotidienne. L'objectif était donc de construire une interface \textbf{professionnelle, intuitive et rassurante}, fidèle à l'image de marque de PwC, tout en rendant lisibles les résultats d'algorithmes d'apprentissage automatique sophistiqués.

Ce chapitre présente les choix technologiques retenus, l'architecture du code, les interfaces utilisateur développées, la manière dont le frontend communique avec les deux backends distincts, ainsi que les mécanismes d'authentification, de contrôle d'accès et de déploiement.

% =====================================================
\section{Stack Technologique}
% =====================================================

\subsection{Framework principal~: Next.js 15 avec App Router}

Le choix de \textbf{Next.js 15} (version la plus récente au moment du développement) avec le système \textbf{App Router} constitue la colonne vertébrale de toute l'interface. Ce framework, construit au-dessus de React 19, offre plusieurs avantages déterminants pour ce type d'application d'entreprise~:

\begin{itemize}[leftmargin=1.5em]
  \item Le \textbf{rendu hybride} permet de combiner rendu côté serveur (pour les données critiques comme les missions et les résultats d'analyse) et rendu côté client (pour les interactions dynamiques comme les filtres et les graphiques), améliorant à la fois la performance perçue et la sécurité.
  \item L'\textbf{App Router} restructure l'organisation du code autour de la notion de layout partagé, de chargement progressif et de routes groupées, ce qui se traduit par une navigation fluide sans rechargement de page.
  \item La gestion intégrée des \textbf{rewrites de proxy} redirige toutes les requêtes vers le backend FastAPI sans exposer l'URL réelle de ce dernier au navigateur, évitant les problèmes CORS et renforçant la sécurité de l'architecture.
\end{itemize}

\textbf{TypeScript 5} est utilisé en mode strict sur l'ensemble du projet. Chaque donnée échangée entre le frontend et les backends est typée explicitement --- missions, utilisateurs, résultats d'analyse, anomalies --- ce qui élimine une classe entière de bugs liés aux données inattendues ou manquantes lors du développement.

\subsection{Stylisation~: TailwindCSS et shadcn/ui}

\textbf{TailwindCSS 3.4} est le système de mise en page et de style. Il adopte une approche utilitaire~: les styles sont composés directement dans le JSX via des classes courtes (\texttt{flex}, \texttt{gap-4}, \texttt{text-lg}, etc.), ce qui accélère le développement tout en garantissant la cohérence visuelle. La palette de couleurs est entièrement configurée pour respecter l'identité PwC~:

\begin{itemize}[leftmargin=1.5em]
  \item \textbf{Orange primaire \texttt{\#D04A02}}~: boutons d'action principaux, accents visuels
  \item \textbf{Bleu foncé \texttt{\#293854}}~: en-têtes et éléments de navigation
  \item \textbf{Niveaux de risque}~: Rouge \texttt{\#C00000} (CRITIQUE), Orange \texttt{\#D04A02} (ÉLEVÉ), Vert \texttt{\#008246} (FAIBLE)
\end{itemize}

\textbf{shadcn/ui} est la bibliothèque de composants retenue. Contrairement à des solutions classiques comme Material UI, shadcn/ui fournit des composants \textit{headless} basés sur Radix UI~: sans style imposé mais avec une accessibilité garantie (navigation clavier, attributs ARIA, compatibilité lecteurs d'écran). Les composants intégrés comprennent Dialog (modales), Tabs (onglets de résultats), Progress (barre de progression d'upload), Select, Dropdown, Badge, Card, Button et Skeleton (placeholders de chargement).

\subsection{Gestion des données~: TanStack React Query 5}

\textbf{TanStack React Query} gère l'ensemble de la récupération, mise en cache et synchronisation des données provenant des APIs. Son fonctionnement est fondamental pour la fluidité de l'expérience utilisateur~:

\begin{itemize}[leftmargin=1.5em]
  \item Les données (missions, utilisateurs, analyses) sont mises en cache \textbf{2 minutes}~: naviguer entre pages ne déclenche pas d'appels réseau répétés.
  \item Après une modification (création de mission, upload de fichier), le cache est \textbf{invalidé et rechargé automatiquement en arrière-plan}, sans que l'utilisateur ait à rafraîchir la page.
  \item Les mutations gèrent automatiquement les états de chargement et d'erreur, déclenchant des notifications via la bibliothèque Sonner.
\end{itemize}

\subsection{Bibliothèques complémentaires}

Le tableau \ref{tab:frontend-libs} récapitule les bibliothèques complémentaires intégrées au projet.

\begin{table}[H]
\centering
\caption{Bibliothèques frontend complémentaires}
\label{tab:frontend-libs}
\begin{tabular}{>{\bfseries}p{4.5cm} p{9.5cm}}
\toprule
Bibliothèque & Usage \\
\midrule
Recharts 2.13 & Graphiques interactifs (camembert, histogramme, barres) basés sur D3.js avec une API React native \\
Lucide React 0.468 & Bibliothèque d'icônes SVG (tree-shaken, seules les icônes utilisées sont embarquées dans le bundle) \\
Sonner 1.7 & Notifications toast positionnées en haut à droite (succès, erreur, information, 4 secondes) \\
React Hook Form 7.54 & Gestion d'état des formulaires sans re-rendu inutile \\
Zod 3.23 & Validation de schéma à l'exécution (email \texttt{@pwc.com}, dates, formats) \\
date-fns 4.1 & Formatage des dates en français (\texttt{dd/MM/yyyy HH:mm}) \\
Axios 1.7 & Client HTTP avec intercepteurs pour l'injection automatique du token JWT \\
Prisma ORM 5.22 & Client base de données avec typage auto-généré pour PostgreSQL \\
bcryptjs 3.0 & Hachage des mots de passe (12 rounds, résistant aux attaques par force brute) \\
jose 5.9 & Génération et vérification des tokens JWT (algorithme HS256) \\
uuid 11.0 & Génération d'identifiants uniques côté client \\
\bottomrule
\end{tabular}
\end{table}

% =====================================================
\section{Architecture du Code}
% =====================================================

L'organisation du code suit les conventions Next.js App Router, structurées en couches logiques clairement séparées~:

\begin{lstlisting}[caption={Structure des dossiers du frontend}, label={lst:frontend-structure}]
frontend/src/
  app/                        # Routes et pages (App Router)
    (auth)/login/             # Page de connexion publique
    (dashboard)/              # Toutes les pages authentifiees
      dashboard/              # Tableau de bord principal
      missions/               # Liste des missions
        [id]/                 # Detail d'une mission
          analysis/           # Assistant d'analyse (wizard 4 etapes)
      history/                # Historique des analyses
      reports/                # Rapports PDF/DOCX generes
      audit-trail/            # Piste d'audit complete
      admin/users/            # Gestion des utilisateurs
    api/                      # Routes API internes Next.js
      auth/                   # Login, logout, me
      missions/               # CRUD missions
      users/                  # Gestion utilisateurs
      analysis-runs/          # Persistance des analyses
      audit-logs/             # Journal d'audit
  components/                 # Composants reutilisables
    analysis/                 # Wizard, tableau de resultats, KPI
    charts/                   # Graphiques Recharts
    datasets/                 # Upload et gestion de fichiers
    missions/                 # Cartes et formulaires de mission
    reports/                  # Generation PDF/DOCX
    explanations/             # Affichage SHAP, LIME, LLM
    layout/                   # Navbar, TopBar
    ui/                       # Primitives shadcn/ui
  lib/
    api/                      # Services d'appel aux APIs
    hooks/                    # Hooks personnalises
    store/                    # Stores locaux (mode demo)
  providers/                  # Providers React (Query, Theme, Language)
  types/                      # Interfaces TypeScript partages
  messages/                   # Traductions (fr.json, en.json)
\end{lstlisting}

% =====================================================
\section{Routes et Interfaces Utilisateur}
% =====================================================

\subsection{Page de Connexion (\texttt{/login})}

La page d'entrée est volontairement sobre et professionnelle. Elle présente le logo PwC et un formulaire à deux champs (email et mot de passe). La validation s'effectue à deux niveaux~:

\begin{itemize}[leftmargin=1.5em]
  \item \textbf{Côté client (Zod)}~: le format email est vérifié et le domaine \texttt{@pwc.com} est obligatoire. Un message d'erreur s'affiche immédiatement sans solliciter le serveur.
  \item \textbf{Côté serveur}~: le mot de passe est comparé au hash bcrypt stocké en base de données (12 rounds, environ 100\,ms par vérification --- délibérément lent pour résister aux attaques par force brute).
\end{itemize}

\subsection{Tableau de Bord Principal (\texttt{/dashboard})}

Le tableau de bord est la première page visible après la connexion. Il offre une vue synthétique de l'activité en cours~: cartes KPI (missions actives, analyses lancées, anomalies détectées), activité récente, et accès rapide aux missions. Le contenu s'adapte automatiquement au \textbf{rôle de l'utilisateur connecté}~: un auditeur ne voit que ses missions assignées, tandis qu'un manager dispose d'une vue d'ensemble complète.

\subsection{Gestion des Missions (\texttt{/missions})}

La liste des missions est présentée sous forme de cartes, chacune affichant le nom de la mission, l'entreprise cliente, le type (commissariat aux comptes, audit interne, due diligence, etc.), les dates, le statut (en préparation, en cours, terminée, archivée) et les auditeurs assignés. Une barre de recherche et des filtres combinés par statut et type permettent de naviguer rapidement.

Le bouton \textit{Nouvelle mission} --- visible uniquement pour les managers, partenaires et administrateurs --- ouvre une modale de création avec un formulaire complet incluant l'affectation d'auditeurs depuis la liste des utilisateurs actifs.

\subsection{Détail d'une Mission (\texttt{/missions/[id]})}

La page de détail consolide~:

\begin{itemize}[leftmargin=1.5em]
  \item Les \textbf{métadonnées de la mission}~: informations générales, équipe, calendrier.
  \item La \textbf{section Datasets}~: liste des fichiers CSV uploadés avec drag-and-drop (\texttt{UploadDropzone}), barre de progression, possibilité de remplacer ou supprimer un fichier (jusqu'à 500\,Mo).
  \item Le \textbf{bouton Lancer l'analyse}~: redirection vers l'assistant multi-étapes.
\end{itemize}

\subsection{Assistant d'Analyse (\texttt{/missions/[id]/analysis})}

C'est l'interface la plus sophistiquée de la plateforme. Elle guide l'utilisateur à travers quatre étapes structurées~:

\begin{description}[leftmargin=1.5em, style=nextline]
  \item[Étape 1 --- Upload du fichier CSV]
    L'utilisateur dépose ou sélectionne le fichier de transactions. Un indicateur de santé du backend ML s'affiche (modèles chargés, LLM disponible) pour rassurer l'utilisateur avant de démarrer.

  \item[Étape 2 --- Sélection du modèle]
    Trois configurations sont proposées avec leur description~:
    \begin{itemize}[leftmargin=1em]
      \item \textbf{Combiné (recommandé)}~: les trois modèles votent en ensemble pour une décision plus robuste.
      \item \textbf{XGBoost uniquement}~: rapide et précis sur les patterns de fraude connus.
      \item \textbf{AutoEncoder uniquement}~: détecte des anomalies inédites, sans supervision.
    \end{itemize}

  \item[Étape 3 --- Exécution de l'analyse]
    Une carte récapitulative s'affiche avant le lancement. Pendant l'analyse, une barre de progression et des messages d'état informent l'utilisateur de chaque phase~: détection du schéma, ingénierie des features, inférence. Le timeout est fixé à 120\,secondes.

  \item[Étape 4 --- Tableau de bord des résultats]
    Les résultats sont organisés en quatre onglets~:
    \begin{itemize}[leftmargin=1em]
      \item \textbf{Graphiques}~: camembert normal/anomalie, histogramme des scores avec ligne de seuil, barres par catégorie.
      \item \textbf{Transactions}~: tableau interactif et filtrable de toutes les transactions avec score de risque et niveau (CRITIQUE/ÉLEVÉ/FAIBLE).
      \item \textbf{Explication}~: pour une transaction sélectionnée, résumé LLM en français, valeurs SHAP, règles LIME, recommandations d'actions.
      \item \textbf{Rapport}~: génération et téléchargement du rapport en PDF ou DOCX.
    \end{itemize}
\end{description}

\subsection{Autres Pages}

\begin{itemize}[leftmargin=1.5em]
  \item \textbf{Historique (\texttt{/history})}~: liste de toutes les analyses précédentes avec filtres, permettant de retrouver et revisualiser des résultats sans relancer une analyse.
  \item \textbf{Rapports (\texttt{/reports})}~: centralise tous les rapports générés, avec téléchargement et archivage.
  \item \textbf{Piste d'audit (\texttt{/audit-trail})}~: journal complet de toutes les actions (connexions, uploads, analyses, rapports) avec horodatage et identité de l'acteur. Réservé aux managers et administrateurs.
  \item \textbf{Gestion des utilisateurs (\texttt{/admin/users})}~: création, modification et désactivation de comptes, assignation de rôles. Réservé aux administrateurs.
\end{itemize}

% =====================================================
\section{Architecture API et Communication avec le Backend}
% =====================================================

\subsection{Double couche d'API}

Le frontend communique avec deux backends distincts selon une architecture en deux niveaux~:

\begin{description}[leftmargin=1.5em, style=nextline]
  \item[Couche 1 --- Next.js API Routes (PostgreSQL)]
    Un backend léger intégré à Next.js, accessible via \texttt{/api/*}, gère toute la persistance métier~: utilisateurs, missions, datasets, historique d'analyses, piste d'audit. Ces routes s'appuient sur Prisma ORM pour interagir avec une base PostgreSQL.

  \item[Couche 2 --- FastAPI ML Backend]
    Le service Python qui héberge les modèles de machine learning, accessible via \texttt{/ml/api/*}. Le frontend ne communique jamais directement avec ce service~: les requêtes transitent par un \textbf{proxy configuré dans \texttt{next.config.ts}}, qui redirige les appels \texttt{/ml/*} vers \texttt{http://localhost:8000}. Ce mécanisme élimine les problèmes CORS et masque l'URL réelle du backend ML au navigateur.
\end{description}

\subsection{Services API}

Le tableau \ref{tab:api-ml} présente les endpoints ML consommés par le frontend, et le tableau \ref{tab:api-app} les endpoints internes.

\begin{table}[H]
\centering
\caption{Endpoints FastAPI ML consommés par le frontend}
\label{tab:api-ml}
\begin{tabular}{p{3.5cm} p{3cm} p{7cm}}
\toprule
\textbf{Service} & \textbf{Méthode / Route} & \textbf{Usage} \\
\midrule
\texttt{predict} & POST \texttt{/api/predict} & Upload CSV, pipeline complet (détection schéma, feature engineering, inférence) \\
\texttt{explain} & GET \texttt{/api/explain/\{tx\_id\}} & SHAP + LIME + explication LLM pour une transaction \\
\texttt{explainBatch} & POST \texttt{/api/explain/batch} & Explications en lot (LIME désactivé pour performance) \\
\texttt{getModels} & GET \texttt{/api/models} & Métriques des modèles (recall, précision, F1, ROC-AUC) \\
\texttt{profile} & POST \texttt{/api/profile} & Profiling qualité du dataset sans prédiction \\
\texttt{health} & GET \texttt{/api/health} & Santé du backend (modèles chargés, LLM disponible) \\
\texttt{generatePDF} & POST \texttt{/api/report} & Génération du rapport PDF \\
\texttt{generateDOCX} & POST \texttt{/api/report/docx} & Génération du rapport Word \\
\bottomrule
\end{tabular}
\end{table}

\begin{table}[H]
\centering
\caption{Endpoints Next.js internes (PostgreSQL)}
\label{tab:api-app}
\begin{tabular}{p{4cm} p{3.5cm} p{6cm}}
\toprule
\textbf{Service} & \textbf{Route} & \textbf{Usage} \\
\midrule
Auth & \texttt{/api/auth/login} & Émission du JWT \\
& \texttt{/api/auth/logout} & Invalidation de session \\
& \texttt{/api/auth/me} & Utilisateur courant \\
Missions & \texttt{/api/missions} & Liste / création \\
& \texttt{/api/missions/[id]} & Lecture / modification / suppression \\
Datasets & \texttt{/api/missions/[id]/datasets} & Liste / upload / remplacement / suppression \\
Utilisateurs & \texttt{/api/users} & Liste et détail \\
Analyses & \texttt{/api/analysis-runs} & Persistance de l'historique \\
Audit & \texttt{/api/audit-logs} & Enregistrement des actions \\
\bottomrule
\end{tabular}
\end{table}

\subsection{Intercepteurs Axios}

Les deux clients HTTP sont configurés avec des intercepteurs~:

\begin{itemize}[leftmargin=1.5em]
  \item \textbf{Intercepteur de requête}~: injecte automatiquement l'en-tête \texttt{Authorization: Bearer \{token\}} à chaque appel, en lisant le token depuis le \texttt{localStorage}.
  \item \textbf{Intercepteur de réponse}~:
    \begin{itemize}[leftmargin=1em]
      \item Réponse \textbf{401} (token expiré ou invalide) → nettoyage du \texttt{localStorage} et redirection vers \texttt{/login}.
      \item \textbf{Erreur réseau} → une tentative de relance automatique.
      \item \textbf{Erreur 4xx} (erreur client) → pas de relance.
    \end{itemize}
\end{itemize}

% =====================================================
\section{Authentification et Contrôle d'Accès}
% =====================================================

\subsection{Flux JWT}

L'authentification repose sur des \textbf{tokens JWT} (JSON Web Tokens) signés avec l'algorithme HS256. Le flux est le suivant~:

\begin{enumerate}[leftmargin=1.5em]
  \item L'utilisateur soumet son email (\texttt{@pwc.com} obligatoire) et son mot de passe.
  \item Le serveur vérifie l'email et compare le mot de passe au hash bcrypt (12 rounds).
  \item Un JWT est généré avec les informations de l'utilisateur (id, email, nom, rôle) et une durée de validité de \textbf{8 heures}.
  \item Le token est retourné au client et stocké dans le \texttt{localStorage} (\texttt{pwc\_token}).
  \item À chaque requête ultérieure, l'intercepteur Axios injecte le token dans l'en-tête \texttt{Authorization}.
  \item Si le token expire, le serveur répond 401 → déconnexion automatique et redirection vers \texttt{/login}.
\end{enumerate}

\subsection{Rôles et Permissions}

Quatre rôles définissent les droits d'accès dans la plateforme~:

\begin{description}[leftmargin=1.5em, style=nextline]
  \item[Auditeur] Accède uniquement aux missions qui lui sont assignées. Peut uploader des fichiers CSV, lancer des analyses et commenter des anomalies.
  \item[Manager] Peut créer et modifier des missions, assigner des auditeurs, approuver des rapports. Accès à toutes les missions.
  \item[Partenaire] Peut créer des missions et consulter les rapports générés.
  \item[Administrateur] Accès complet à toutes les fonctionnalités, y compris la création et la modification des comptes utilisateurs.
\end{description}

Le hook personnalisé \texttt{usePermissions()} est utilisé dans chaque composant pour afficher ou masquer conditionnellement les éléments d'interface selon le rôle de l'utilisateur connecté.

% =====================================================
\section{Déploiement des Modèles côté Frontend}
% =====================================================

\subsection{Sélection des modèles}

L'étape 2 de l'assistant d'analyse (\texttt{AnalysisWizard.tsx}) présente trois configurations de modèles. Le choix de l'utilisateur est transmis en paramètre (\texttt{model\_type}) dans le corps de la requête à \texttt{POST /ml/api/predict}, avec les valeurs possibles \texttt{combined}, \texttt{xgboost} ou \texttt{autoencoder}.

\subsection{Orchestration de l'analyse}

Quand l'utilisateur déclenche une analyse~:

\begin{enumerate}[leftmargin=1.5em]
  \item Le fichier CSV est envoyé en \textit{multipart/form-data} avec le type de modèle choisi.
  \item Une barre de progression s'affiche --- le timeout est fixé à 120 secondes pour accommoder les grands fichiers.
  \item Le backend FastAPI exécute séquentiellement~: détection du schéma, profiling qualité, ingénierie des 14 features dérivées, puis inférence des modèles sélectionnés.
  \item La réponse contient l'ensemble des transactions avec leurs scores et niveaux de risque.
  \item Si la détection de schéma est incomplète (confiance sous le seuil), un \textbf{avertissement} est affiché pour informer l'auditeur de l'impact possible sur la précision.
\end{enumerate}

\subsection{Explications à la demande}

Les explications (SHAP, LIME, résumé LLM) ne sont pas calculées en masse au moment de l'analyse~: elles sont générées \textbf{à la demande}, pour la transaction que l'auditeur sélectionne dans le tableau. Cela évite une surcharge inutile du backend pour des transactions qui ne seront jamais inspectées en détail. Pour la génération de rapport, un appel batch (\texttt{POST /ml/api/explain/batch}) permet de récupérer les explications des anomalies principales en une seule requête, avec LIME désactivé pour des raisons de performance.

% =====================================================
\section{Internationalisation et Thème}
% =====================================================

\subsection{Interface bilingue français/anglais}

L'interface est entièrement bilingue. Les traductions sont stockées dans \texttt{src/messages/fr.json} et \texttt{en.json}. La langue est persistée dans le \texttt{localStorage} (\texttt{pwc\_locale}) et peut être changée à tout moment via un sélecteur dans la barre de navigation. Le \texttt{LanguageContext} injecte la fonction de traduction \texttt{t()} dans tous les composants qui en ont besoin, avec support de l'interpolation de variables.

\subsection{Mode sombre}

Un bouton dans la barre de navigation permet de basculer entre mode clair et mode sombre. Le \texttt{ThemeProvider} gère cela via la classe \texttt{dark} sur l'élément racine HTML. TailwindCSS applique automatiquement les variantes \texttt{dark:} définies sur chaque composant.

% =====================================================
\section{Performance et Configuration}
% =====================================================

\subsection{Optimisations de build (\texttt{next.config.ts})}

\begin{itemize}[leftmargin=1.5em]
  \item \textbf{Tree-shaking des imports baril}~: seules les icônes et composants réellement utilisés de \texttt{lucide-react}, \texttt{recharts} et \texttt{@radix-ui} sont inclus dans le bundle final.
  \item \textbf{Cache agressif des assets statiques}~: 31\,536\,000 secondes (1 an) pour les images et les fichiers Next.js compilés.
  \item \textbf{Optimisation des images}~: conversion automatique en AVIF et WebP (40 à 50\,\% plus légers que JPEG/PNG).
  \item \textbf{Compression}~: activée sur toutes les réponses HTTP.
  \item \textbf{Standalone output}~: pour Docker, le build n'embarque que les dépendances strictement nécessaires à l'exécution, produisant une image légère.
  \item \textbf{Limite de 500\,Mo}~: configurée pour les Server Actions afin de permettre l'upload de fichiers CSV volumineux.
\end{itemize}

\subsection{Schéma de base de données}

Le schéma Prisma (\texttt{prisma/schema.prisma}) définit 11 modèles interconnectés~: \texttt{User}, \texttt{Mission}, \texttt{MissionAssignment}, \texttt{Dataset}, \texttt{DatasetVersion}, \texttt{AnalysisRun}, \texttt{Anomaly}, \texttt{AnomalyComment}, \texttt{Report}, \texttt{AuditLog} et \texttt{Permission}. Des index sont définis sur les colonnes les plus fréquemment filtrées (\texttt{userId}, \texttt{missionId}, \texttt{action}, \texttt{timestamp}) pour maintenir les performances de requête à mesure que les données s'accumulent.

% =====================================================
\section{Gestion des Erreurs et Expérience Utilisateur}
% =====================================================

Chaque opération susceptible d'échouer est traitée explicitement~:

\begin{itemize}[leftmargin=1.5em]
  \item \textbf{Upload échoué}~: message descriptif formaté depuis l'erreur serveur, avec possibilité de réessayer.
  \item \textbf{Analyse échouée}~: message extrait de la réponse FastAPI, avec code d'erreur si disponible.
  \item \textbf{Token expiré}~: déconnexion propre et redirection sans page blanche ni message cryptique.
  \item \textbf{Backend indisponible}~: indicateur rouge dans l'assistant d'analyse avec message d'alerte explicite.
  \item \textbf{Schéma CSV non reconnu}~: avertissement détaillé sur les colonnes manquantes et leur impact potentiel sur la précision des modèles.
\end{itemize}

Les \textbf{notifications Sonner} (4 secondes, en haut à droite) confirment toutes les actions réussies~: mission créée, fichier uploadé, analyse terminée, rapport généré.

Les \textbf{états de chargement} sont gérés avec des squelettes (\texttt{Skeleton}) qui maintiennent la mise en page pendant les requêtes, évitant les sauts visuels qui nuisent à la perception de fiabilité de l'outil.

% =====================================================
\section*{Conclusion du chapitre}
\addcontentsline{toc}{section}{Conclusion du chapitre}
% =====================================================

Le frontend de cette plateforme est une application React de niveau professionnel, pensée pour un contexte d'entreprise exigeant. Sa cohérence repose sur une architecture claire~: Next.js 15 pour le rendu hybride et le proxy ML, TanStack Query pour la gestion de l'état serveur, shadcn/ui pour l'accessibilité, et Axios avec intercepteurs pour la sécurité des échanges.

La force de cette interface réside dans la \textbf{lisibilité du parcours utilisateur}~: du premier clic sur \textit{Connexion} jusqu'au téléchargement du rapport PDF, chaque étape est guidée et protégée par des validations, des messages d'erreur clairs et des indicateurs de progression. La complexité des modèles d'intelligence artificielle est totalement abstraite~: l'auditeur choisit un fichier, sélectionne une configuration, lance l'analyse, et reçoit des résultats exploitables en quelques dizaines de secondes, accompagnés d'explications rédigées en français naturel.

% =====================================================

\chapter{Interface Frontend de la Plateforme}

\section*{Introduction du chapitre}
\addcontentsline{toc}{section}{Introduction du chapitre}

La conception de l'interface frontend a répondu à un défi central~: mettre entre les mains d'auditeurs, dont la majorité ne sont pas informaticiens, un outil d'intelligence artificielle puissant, sans que la complexité technique ne transparaisse dans leur expérience quotidienne. L'objectif était donc de construire une interface \textbf{professionnelle, intuitive et rassurante}, fidèle à l'image de marque de PwC, tout en rendant lisibles les résultats d'algorithmes d'apprentissage automatique sophistiqués.

Ce chapitre présente les choix technologiques retenus, l'architecture du code, les interfaces utilisateur développées, la manière dont le frontend communique avec les deux backends distincts, ainsi que les mécanismes d'authentification, de contrôle d'accès et de déploiement.

% =====================================================
\section{Stack Technologique}
% =====================================================

\subsection{Framework principal~: Next.js 15 avec App Router}

Le choix de \textbf{Next.js 15} (version la plus récente au moment du développement) avec le système \textbf{App Router} constitue la colonne vertébrale de toute l'interface. Ce framework, construit au-dessus de React 19, offre plusieurs avantages déterminants pour ce type d'application d'entreprise~:

\begin{itemize}[leftmargin=1.5em]
  \item Le \textbf{rendu hybride} permet de combiner rendu côté serveur (pour les données critiques comme les missions et les résultats d'analyse) et rendu côté client (pour les interactions dynamiques comme les filtres et les graphiques), améliorant à la fois la performance perçue et la sécurité.
  \item L'\textbf{App Router} restructure l'organisation du code autour de la notion de layout partagé, de chargement progressif et de routes groupées, ce qui se traduit par une navigation fluide sans rechargement de page.
  \item La gestion intégrée des \textbf{rewrites de proxy} redirige toutes les requêtes vers le backend FastAPI sans exposer l'URL réelle de ce dernier au navigateur, évitant les problèmes CORS et renforçant la sécurité de l'architecture.
\end{itemize}

\textbf{TypeScript 5} est utilisé en mode strict sur l'ensemble du projet. Chaque donnée échangée entre le frontend et les backends est typée explicitement --- missions, utilisateurs, résultats d'analyse, anomalies --- ce qui élimine une classe entière de bugs liés aux données inattendues ou manquantes lors du développement.

\subsection{Stylisation~: TailwindCSS et shadcn/ui}

\textbf{TailwindCSS 3.4} est le système de mise en page et de style. Il adopte une approche utilitaire~: les styles sont composés directement dans le JSX via des classes courtes (\texttt{flex}, \texttt{gap-4}, \texttt{text-lg}, etc.), ce qui accélère le développement tout en garantissant la cohérence visuelle. La palette de couleurs est entièrement configurée pour respecter l'identité PwC~:

\begin{itemize}[leftmargin=1.5em]
  \item \textbf{Orange primaire \texttt{\#D04A02}}~: boutons d'action principaux, accents visuels
  \item \textbf{Bleu foncé \texttt{\#293854}}~: en-têtes et éléments de navigation
  \item \textbf{Niveaux de risque}~: Rouge \texttt{\#C00000} (CRITIQUE), Orange \texttt{\#D04A02} (ÉLEVÉ), Vert \texttt{\#008246} (FAIBLE)
\end{itemize}

\textbf{shadcn/ui} est la bibliothèque de composants retenue. Contrairement à des solutions classiques comme Material UI, shadcn/ui fournit des composants \textit{headless} basés sur Radix UI~: sans style imposé mais avec une accessibilité garantie (navigation clavier, attributs ARIA, compatibilité lecteurs d'écran). Les composants intégrés comprennent Dialog (modales), Tabs (onglets de résultats), Progress (barre de progression d'upload), Select, Dropdown, Badge, Card, Button et Skeleton (placeholders de chargement).

\subsection{Gestion des données~: TanStack React Query 5}

\textbf{TanStack React Query} gère l'ensemble de la récupération, mise en cache et synchronisation des données provenant des APIs. Son fonctionnement est fondamental pour la fluidité de l'expérience utilisateur~:

\begin{itemize}[leftmargin=1.5em]
  \item Les données (missions, utilisateurs, analyses) sont mises en cache \textbf{2 minutes}~: naviguer entre pages ne déclenche pas d'appels réseau répétés.
  \item Après une modification (création de mission, upload de fichier), le cache est \textbf{invalidé et rechargé automatiquement en arrière-plan}, sans que l'utilisateur ait à rafraîchir la page.
  \item Les mutations gèrent automatiquement les états de chargement et d'erreur, déclenchant des notifications via la bibliothèque Sonner.
\end{itemize}

\subsection{Bibliothèques complémentaires}

Le tableau \ref{tab:frontend-libs} récapitule les bibliothèques complémentaires intégrées au projet.

\begin{table}[H]
\centering
\caption{Bibliothèques frontend complémentaires}
\label{tab:frontend-libs}
\begin{tabular}{>{\bfseries}p{4.5cm} p{9.5cm}}
\toprule
Bibliothèque & Usage \\
\midrule
Recharts 2.13 & Graphiques interactifs (camembert, histogramme, barres) basés sur D3.js avec une API React native \\
Lucide React 0.468 & Bibliothèque d'icônes SVG (tree-shaken, seules les icônes utilisées sont embarquées dans le bundle) \\
Sonner 1.7 & Notifications toast positionnées en haut à droite (succès, erreur, information, 4 secondes) \\
React Hook Form 7.54 & Gestion d'état des formulaires sans re-rendu inutile \\
Zod 3.23 & Validation de schéma à l'exécution (email \texttt{@pwc.com}, dates, formats) \\
date-fns 4.1 & Formatage des dates en français (\texttt{dd/MM/yyyy HH:mm}) \\
Axios 1.7 & Client HTTP avec intercepteurs pour l'injection automatique du token JWT \\
Prisma ORM 5.22 & Client base de données avec typage auto-généré pour PostgreSQL \\
bcryptjs 3.0 & Hachage des mots de passe (12 rounds, résistant aux attaques par force brute) \\
jose 5.9 & Génération et vérification des tokens JWT (algorithme HS256) \\
uuid 11.0 & Génération d'identifiants uniques côté client \\
\bottomrule
\end{tabular}
\end{table}

% =====================================================
\section{Architecture du Code}
% =====================================================

L'organisation du code suit les conventions Next.js App Router, structurées en couches logiques clairement séparées~:

\begin{lstlisting}[caption={Structure des dossiers du frontend}, label={lst:frontend-structure}]
frontend/src/
  app/                        # Routes et pages (App Router)
    (auth)/login/             # Page de connexion publique
    (dashboard)/              # Toutes les pages authentifiees
      dashboard/              # Tableau de bord principal
      missions/               # Liste des missions
        [id]/                 # Detail d'une mission
          analysis/           # Assistant d'analyse (wizard 4 etapes)
      history/                # Historique des analyses
      reports/                # Rapports PDF/DOCX generes
      audit-trail/            # Piste d'audit complete
      admin/users/            # Gestion des utilisateurs
    api/                      # Routes API internes Next.js
      auth/                   # Login, logout, me
      missions/               # CRUD missions
      users/                  # Gestion utilisateurs
      analysis-runs/          # Persistance des analyses
      audit-logs/             # Journal d'audit
  components/                 # Composants reutilisables
    analysis/                 # Wizard, tableau de resultats, KPI
    charts/                   # Graphiques Recharts
    datasets/                 # Upload et gestion de fichiers
    missions/                 # Cartes et formulaires de mission
    reports/                  # Generation PDF/DOCX
    explanations/             # Affichage SHAP, LIME, LLM
    layout/                   # Navbar, TopBar
    ui/                       # Primitives shadcn/ui
  lib/
    api/                      # Services d'appel aux APIs
    hooks/                    # Hooks personnalises
    store/                    # Stores locaux (mode demo)
  providers/                  # Providers React (Query, Theme, Language)
  types/                      # Interfaces TypeScript partages
  messages/                   # Traductions (fr.json, en.json)
\end{lstlisting}

% =====================================================
\section{Routes et Interfaces Utilisateur}
% =====================================================

\subsection{Page de Connexion (\texttt{/login})}

La page d'entrée est volontairement sobre et professionnelle. Elle présente le logo PwC et un formulaire à deux champs (email et mot de passe). La validation s'effectue à deux niveaux~:

\begin{itemize}[leftmargin=1.5em]
  \item \textbf{Côté client (Zod)}~: le format email est vérifié et le domaine \texttt{@pwc.com} est obligatoire. Un message d'erreur s'affiche immédiatement sans solliciter le serveur.
  \item \textbf{Côté serveur}~: le mot de passe est comparé au hash bcrypt stocké en base de données (12 rounds, environ 100\,ms par vérification --- délibérément lent pour résister aux attaques par force brute).
\end{itemize}

\subsection{Tableau de Bord Principal (\texttt{/dashboard})}

Le tableau de bord est la première page visible après la connexion. Il offre une vue synthétique de l'activité en cours~: cartes KPI (missions actives, analyses lancées, anomalies détectées), activité récente, et accès rapide aux missions. Le contenu s'adapte automatiquement au \textbf{rôle de l'utilisateur connecté}~: un auditeur ne voit que ses missions assignées, tandis qu'un manager dispose d'une vue d'ensemble complète.

\subsection{Gestion des Missions (\texttt{/missions})}

La liste des missions est présentée sous forme de cartes, chacune affichant le nom de la mission, l'entreprise cliente, le type (commissariat aux comptes, audit interne, due diligence, etc.), les dates, le statut (en préparation, en cours, terminée, archivée) et les auditeurs assignés. Une barre de recherche et des filtres combinés par statut et type permettent de naviguer rapidement.

Le bouton \textit{Nouvelle mission} --- visible uniquement pour les managers, partenaires et administrateurs --- ouvre une modale de création avec un formulaire complet incluant l'affectation d'auditeurs depuis la liste des utilisateurs actifs.

\subsection{Détail d'une Mission (\texttt{/missions/[id]})}

La page de détail consolide~:

\begin{itemize}[leftmargin=1.5em]
  \item Les \textbf{métadonnées de la mission}~: informations générales, équipe, calendrier.
  \item La \textbf{section Datasets}~: liste des fichiers CSV uploadés avec drag-and-drop (\texttt{UploadDropzone}), barre de progression, possibilité de remplacer ou supprimer un fichier (jusqu'à 500\,Mo).
  \item Le \textbf{bouton Lancer l'analyse}~: redirection vers l'assistant multi-étapes.
\end{itemize}

\subsection{Assistant d'Analyse (\texttt{/missions/[id]/analysis})}

C'est l'interface la plus sophistiquée de la plateforme. Elle guide l'utilisateur à travers quatre étapes structurées~:

\begin{description}[leftmargin=1.5em, style=nextline]
  \item[Étape 1 --- Upload du fichier CSV]
    L'utilisateur dépose ou sélectionne le fichier de transactions. Un indicateur de santé du backend ML s'affiche (modèles chargés, LLM disponible) pour rassurer l'utilisateur avant de démarrer.

  \item[Étape 2 --- Sélection du modèle]
    Trois configurations sont proposées avec leur description~:
    \begin{itemize}[leftmargin=1em]
      \item \textbf{Combiné (recommandé)}~: les trois modèles votent en ensemble pour une décision plus robuste.
      \item \textbf{XGBoost uniquement}~: rapide et précis sur les patterns de fraude connus.
      \item \textbf{AutoEncoder uniquement}~: détecte des anomalies inédites, sans supervision.
    \end{itemize}

  \item[Étape 3 --- Exécution de l'analyse]
    Une carte récapitulative s'affiche avant le lancement. Pendant l'analyse, une barre de progression et des messages d'état informent l'utilisateur de chaque phase~: détection du schéma, ingénierie des features, inférence. Le timeout est fixé à 120\,secondes.

  \item[Étape 4 --- Tableau de bord des résultats]
    Les résultats sont organisés en quatre onglets~:
    \begin{itemize}[leftmargin=1em]
      \item \textbf{Graphiques}~: camembert normal/anomalie, histogramme des scores avec ligne de seuil, barres par catégorie.
      \item \textbf{Transactions}~: tableau interactif et filtrable de toutes les transactions avec score de risque et niveau (CRITIQUE/ÉLEVÉ/FAIBLE).
      \item \textbf{Explication}~: pour une transaction sélectionnée, résumé LLM en français, valeurs SHAP, règles LIME, recommandations d'actions.
      \item \textbf{Rapport}~: génération et téléchargement du rapport en PDF ou DOCX.
    \end{itemize}
\end{description}

\subsection{Autres Pages}

\begin{itemize}[leftmargin=1.5em]
  \item \textbf{Historique (\texttt{/history})}~: liste de toutes les analyses précédentes avec filtres, permettant de retrouver et revisualiser des résultats sans relancer une analyse.
  \item \textbf{Rapports (\texttt{/reports})}~: centralise tous les rapports générés, avec téléchargement et archivage.
  \item \textbf{Piste d'audit (\texttt{/audit-trail})}~: journal complet de toutes les actions (connexions, uploads, analyses, rapports) avec horodatage et identité de l'acteur. Réservé aux managers et administrateurs.
  \item \textbf{Gestion des utilisateurs (\texttt{/admin/users})}~: création, modification et désactivation de comptes, assignation de rôles. Réservé aux administrateurs.
\end{itemize}

% =====================================================
\section{Architecture API et Communication avec le Backend}
% =====================================================

\subsection{Double couche d'API}

Le frontend communique avec deux backends distincts selon une architecture en deux niveaux~:

\begin{description}[leftmargin=1.5em, style=nextline]
  \item[Couche 1 --- Next.js API Routes (PostgreSQL)]
    Un backend léger intégré à Next.js, accessible via \texttt{/api/*}, gère toute la persistance métier~: utilisateurs, missions, datasets, historique d'analyses, piste d'audit. Ces routes s'appuient sur Prisma ORM pour interagir avec une base PostgreSQL.

  \item[Couche 2 --- FastAPI ML Backend]
    Le service Python qui héberge les modèles de machine learning, accessible via \texttt{/ml/api/*}. Le frontend ne communique jamais directement avec ce service~: les requêtes transitent par un \textbf{proxy configuré dans \texttt{next.config.ts}}, qui redirige les appels \texttt{/ml/*} vers \texttt{http://localhost:8000}. Ce mécanisme élimine les problèmes CORS et masque l'URL réelle du backend ML au navigateur.
\end{description}

\subsection{Services API}

Le tableau \ref{tab:api-ml} présente les endpoints ML consommés par le frontend, et le tableau \ref{tab:api-app} les endpoints internes.

\begin{table}[H]
\centering
\caption{Endpoints FastAPI ML consommés par le frontend}
\label{tab:api-ml}
\begin{tabular}{p{3.5cm} p{3cm} p{7cm}}
\toprule
\textbf{Service} & \textbf{Méthode / Route} & \textbf{Usage} \\
\midrule
\texttt{predict} & POST \texttt{/api/predict} & Upload CSV, pipeline complet (détection schéma, feature engineering, inférence) \\
\texttt{explain} & GET \texttt{/api/explain/\{tx\_id\}} & SHAP + LIME + explication LLM pour une transaction \\
\texttt{explainBatch} & POST \texttt{/api/explain/batch} & Explications en lot (LIME désactivé pour performance) \\
\texttt{getModels} & GET \texttt{/api/models} & Métriques des modèles (recall, précision, F1, ROC-AUC) \\
\texttt{profile} & POST \texttt{/api/profile} & Profiling qualité du dataset sans prédiction \\
\texttt{health} & GET \texttt{/api/health} & Santé du backend (modèles chargés, LLM disponible) \\
\texttt{generatePDF} & POST \texttt{/api/report} & Génération du rapport PDF \\
\texttt{generateDOCX} & POST \texttt{/api/report/docx} & Génération du rapport Word \\
\bottomrule
\end{tabular}
\end{table}

\begin{table}[H]
\centering
\caption{Endpoints Next.js internes (PostgreSQL)}
\label{tab:api-app}
\begin{tabular}{p{4cm} p{3.5cm} p{6cm}}
\toprule
\textbf{Service} & \textbf{Route} & \textbf{Usage} \\
\midrule
Auth & \texttt{/api/auth/login} & Émission du JWT \\
& \texttt{/api/auth/logout} & Invalidation de session \\
& \texttt{/api/auth/me} & Utilisateur courant \\
Missions & \texttt{/api/missions} & Liste / création \\
& \texttt{/api/missions/[id]} & Lecture / modification / suppression \\
Datasets & \texttt{/api/missions/[id]/datasets} & Liste / upload / remplacement / suppression \\
Utilisateurs & \texttt{/api/users} & Liste et détail \\
Analyses & \texttt{/api/analysis-runs} & Persistance de l'historique \\
Audit & \texttt{/api/audit-logs} & Enregistrement des actions \\
\bottomrule
\end{tabular}
\end{table}

\subsection{Intercepteurs Axios}

Les deux clients HTTP sont configurés avec des intercepteurs~:

\begin{itemize}[leftmargin=1.5em]
  \item \textbf{Intercepteur de requête}~: injecte automatiquement l'en-tête \texttt{Authorization: Bearer \{token\}} à chaque appel, en lisant le token depuis le \texttt{localStorage}.
  \item \textbf{Intercepteur de réponse}~:
    \begin{itemize}[leftmargin=1em]
      \item Réponse \textbf{401} (token expiré ou invalide) → nettoyage du \texttt{localStorage} et redirection vers \texttt{/login}.
      \item \textbf{Erreur réseau} → une tentative de relance automatique.
      \item \textbf{Erreur 4xx} (erreur client) → pas de relance.
    \end{itemize}
\end{itemize}

% =====================================================
\section{Authentification et Contrôle d'Accès}
% =====================================================

\subsection{Flux JWT}

L'authentification repose sur des \textbf{tokens JWT} (JSON Web Tokens) signés avec l'algorithme HS256. Le flux est le suivant~:

\begin{enumerate}[leftmargin=1.5em]
  \item L'utilisateur soumet son email (\texttt{@pwc.com} obligatoire) et son mot de passe.
  \item Le serveur vérifie l'email et compare le mot de passe au hash bcrypt (12 rounds).
  \item Un JWT est généré avec les informations de l'utilisateur (id, email, nom, rôle) et une durée de validité de \textbf{8 heures}.
  \item Le token est retourné au client et stocké dans le \texttt{localStorage} (\texttt{pwc\_token}).
  \item À chaque requête ultérieure, l'intercepteur Axios injecte le token dans l'en-tête \texttt{Authorization}.
  \item Si le token expire, le serveur répond 401 → déconnexion automatique et redirection vers \texttt{/login}.
\end{enumerate}

\subsection{Rôles et Permissions}

Quatre rôles définissent les droits d'accès dans la plateforme~:

\begin{description}[leftmargin=1.5em, style=nextline]
  \item[Auditeur] Accède uniquement aux missions qui lui sont assignées. Peut uploader des fichiers CSV, lancer des analyses et commenter des anomalies.
  \item[Manager] Peut créer et modifier des missions, assigner des auditeurs, approuver des rapports. Accès à toutes les missions.
  \item[Partenaire] Peut créer des missions et consulter les rapports générés.
  \item[Administrateur] Accès complet à toutes les fonctionnalités, y compris la création et la modification des comptes utilisateurs.
\end{description}

Le hook personnalisé \texttt{usePermissions()} est utilisé dans chaque composant pour afficher ou masquer conditionnellement les éléments d'interface selon le rôle de l'utilisateur connecté.

% =====================================================
\section{Déploiement des Modèles côté Frontend}
% =====================================================

\subsection{Sélection des modèles}

L'étape 2 de l'assistant d'analyse (\texttt{AnalysisWizard.tsx}) présente trois configurations de modèles. Le choix de l'utilisateur est transmis en paramètre (\texttt{model\_type}) dans le corps de la requête à \texttt{POST /ml/api/predict}, avec les valeurs possibles \texttt{combined}, \texttt{xgboost} ou \texttt{autoencoder}.

\subsection{Orchestration de l'analyse}

Quand l'utilisateur déclenche une analyse~:

\begin{enumerate}[leftmargin=1.5em]
  \item Le fichier CSV est envoyé en \textit{multipart/form-data} avec le type de modèle choisi.
  \item Une barre de progression s'affiche --- le timeout est fixé à 120 secondes pour accommoder les grands fichiers.
  \item Le backend FastAPI exécute séquentiellement~: détection du schéma, profiling qualité, ingénierie des 14 features dérivées, puis inférence des modèles sélectionnés.
  \item La réponse contient l'ensemble des transactions avec leurs scores et niveaux de risque.
  \item Si la détection de schéma est incomplète (confiance sous le seuil), un \textbf{avertissement} est affiché pour informer l'auditeur de l'impact possible sur la précision.
\end{enumerate}

\subsection{Explications à la demande}

Les explications (SHAP, LIME, résumé LLM) ne sont pas calculées en masse au moment de l'analyse~: elles sont générées \textbf{à la demande}, pour la transaction que l'auditeur sélectionne dans le tableau. Cela évite une surcharge inutile du backend pour des transactions qui ne seront jamais inspectées en détail. Pour la génération de rapport, un appel batch (\texttt{POST /ml/api/explain/batch}) permet de récupérer les explications des anomalies principales en une seule requête, avec LIME désactivé pour des raisons de performance.

% =====================================================
\section{Internationalisation et Thème}
% =====================================================

\subsection{Interface bilingue français/anglais}

L'interface est entièrement bilingue. Les traductions sont stockées dans \texttt{src/messages/fr.json} et \texttt{en.json}. La langue est persistée dans le \texttt{localStorage} (\texttt{pwc\_locale}) et peut être changée à tout moment via un sélecteur dans la barre de navigation. Le \texttt{LanguageContext} injecte la fonction de traduction \texttt{t()} dans tous les composants qui en ont besoin, avec support de l'interpolation de variables.

\subsection{Mode sombre}

Un bouton dans la barre de navigation permet de basculer entre mode clair et mode sombre. Le \texttt{ThemeProvider} gère cela via la classe \texttt{dark} sur l'élément racine HTML. TailwindCSS applique automatiquement les variantes \texttt{dark:} définies sur chaque composant.

% =====================================================
\section{Performance et Configuration}
% =====================================================

\subsection{Optimisations de build (\texttt{next.config.ts})}

\begin{itemize}[leftmargin=1.5em]
  \item \textbf{Tree-shaking des imports baril}~: seules les icônes et composants réellement utilisés de \texttt{lucide-react}, \texttt{recharts} et \texttt{@radix-ui} sont inclus dans le bundle final.
  \item \textbf{Cache agressif des assets statiques}~: 31\,536\,000 secondes (1 an) pour les images et les fichiers Next.js compilés.
  \item \textbf{Optimisation des images}~: conversion automatique en AVIF et WebP (40 à 50\,\% plus légers que JPEG/PNG).
  \item \textbf{Compression}~: activée sur toutes les réponses HTTP.
  \item \textbf{Standalone output}~: pour Docker, le build n'embarque que les dépendances strictement nécessaires à l'exécution, produisant une image légère.
  \item \textbf{Limite de 500\,Mo}~: configurée pour les Server Actions afin de permettre l'upload de fichiers CSV volumineux.
\end{itemize}

\subsection{Schéma de base de données}

Le schéma Prisma (\texttt{prisma/schema.prisma}) définit 11 modèles interconnectés~: \texttt{User}, \texttt{Mission}, \texttt{MissionAssignment}, \texttt{Dataset}, \texttt{DatasetVersion}, \texttt{AnalysisRun}, \texttt{Anomaly}, \texttt{AnomalyComment}, \texttt{Report}, \texttt{AuditLog} et \texttt{Permission}. Des index sont définis sur les colonnes les plus fréquemment filtrées (\texttt{userId}, \texttt{missionId}, \texttt{action}, \texttt{timestamp}) pour maintenir les performances de requête à mesure que les données s'accumulent.

% =====================================================
\section{Gestion des Erreurs et Expérience Utilisateur}
% =====================================================

Chaque opération susceptible d'échouer est traitée explicitement~:

\begin{itemize}[leftmargin=1.5em]
  \item \textbf{Upload échoué}~: message descriptif formaté depuis l'erreur serveur, avec possibilité de réessayer.
  \item \textbf{Analyse échouée}~: message extrait de la réponse FastAPI, avec code d'erreur si disponible.
  \item \textbf{Token expiré}~: déconnexion propre et redirection sans page blanche ni message cryptique.
  \item \textbf{Backend indisponible}~: indicateur rouge dans l'assistant d'analyse avec message d'alerte explicite.
  \item \textbf{Schéma CSV non reconnu}~: avertissement détaillé sur les colonnes manquantes et leur impact potentiel sur la précision des modèles.
\end{itemize}

Les \textbf{notifications Sonner} (4 secondes, en haut à droite) confirment toutes les actions réussies~: mission créée, fichier uploadé, analyse terminée, rapport généré.

Les \textbf{états de chargement} sont gérés avec des squelettes (\texttt{Skeleton}) qui maintiennent la mise en page pendant les requêtes, évitant les sauts visuels qui nuisent à la perception de fiabilité de l'outil.

% =====================================================
\section*{Conclusion du chapitre}
\addcontentsline{toc}{section}{Conclusion du chapitre}
% =====================================================

Le frontend de cette plateforme est une application React de niveau professionnel, pensée pour un contexte d'entreprise exigeant. Sa cohérence repose sur une architecture claire~: Next.js 15 pour le rendu hybride et le proxy ML, TanStack Query pour la gestion de l'état serveur, shadcn/ui pour l'accessibilité, et Axios avec intercepteurs pour la sécurité des échanges.

La force de cette interface réside dans la \textbf{lisibilité du parcours utilisateur}~: du premier clic sur \textit{Connexion} jusqu'au téléchargement du rapport PDF, chaque étape est guidée et protégée par des validations, des messages d'erreur clairs et des indicateurs de progression. La complexité des modèles d'intelligence artificielle est totalement abstraite~: l'auditeur choisit un fichier, sélectionne une configuration, lance l'analyse, et reçoit des résultats exploitables en quelques dizaines de secondes, accompagnés d'explications rédigées en français naturel.

% =====================================================

\chapter{Interface Frontend de la Plateforme}

\section*{Introduction du chapitre}
\addcontentsline{toc}{section}{Introduction du chapitre}

La conception de l'interface frontend a répondu à un défi central~: mettre entre les mains d'auditeurs, dont la majorité ne sont pas informaticiens, un outil d'intelligence artificielle puissant, sans que la complexité technique ne transparaisse dans leur expérience quotidienne. L'objectif était donc de construire une interface \textbf{professionnelle, intuitive et rassurante}, fidèle à l'image de marque de PwC, tout en rendant lisibles les résultats d'algorithmes d'apprentissage automatique sophistiqués.

Ce chapitre présente les choix technologiques retenus, l'architecture du code, les interfaces utilisateur développées, la manière dont le frontend communique avec les deux backends distincts, ainsi que les mécanismes d'authentification, de contrôle d'accès et de déploiement.

% =====================================================
\section{Stack Technologique}
% =====================================================

\subsection{Framework principal~: Next.js 15 avec App Router}

Le choix de \textbf{Next.js 15} (version la plus récente au moment du développement) avec le système \textbf{App Router} constitue la colonne vertébrale de toute l'interface. Ce framework, construit au-dessus de React 19, offre plusieurs avantages déterminants pour ce type d'application d'entreprise~:

\begin{itemize}[leftmargin=1.5em]
  \item Le \textbf{rendu hybride} permet de combiner rendu côté serveur (pour les données critiques comme les missions et les résultats d'analyse) et rendu côté client (pour les interactions dynamiques comme les filtres et les graphiques), améliorant à la fois la performance perçue et la sécurité.
  \item L'\textbf{App Router} restructure l'organisation du code autour de la notion de layout partagé, de chargement progressif et de routes groupées, ce qui se traduit par une navigation fluide sans rechargement de page.
  \item La gestion intégrée des \textbf{rewrites de proxy} redirige toutes les requêtes vers le backend FastAPI sans exposer l'URL réelle de ce dernier au navigateur, évitant les problèmes CORS et renforçant la sécurité de l'architecture.
\end{itemize}

\textbf{TypeScript 5} est utilisé en mode strict sur l'ensemble du projet. Chaque donnée échangée entre le frontend et les backends est typée explicitement --- missions, utilisateurs, résultats d'analyse, anomalies --- ce qui élimine une classe entière de bugs liés aux données inattendues ou manquantes lors du développement.

\subsection{Stylisation~: TailwindCSS et shadcn/ui}

\textbf{TailwindCSS 3.4} est le système de mise en page et de style. Il adopte une approche utilitaire~: les styles sont composés directement dans le JSX via des classes courtes (\texttt{flex}, \texttt{gap-4}, \texttt{text-lg}, etc.), ce qui accélère le développement tout en garantissant la cohérence visuelle. La palette de couleurs est entièrement configurée pour respecter l'identité PwC~:

\begin{itemize}[leftmargin=1.5em]
  \item \textbf{Orange primaire \texttt{\#D04A02}}~: boutons d'action principaux, accents visuels
  \item \textbf{Bleu foncé \texttt{\#293854}}~: en-têtes et éléments de navigation
  \item \textbf{Niveaux de risque}~: Rouge \texttt{\#C00000} (CRITIQUE), Orange \texttt{\#D04A02} (ÉLEVÉ), Vert \texttt{\#008246} (FAIBLE)
\end{itemize}

\textbf{shadcn/ui} est la bibliothèque de composants retenue. Contrairement à des solutions classiques comme Material UI, shadcn/ui fournit des composants \textit{headless} basés sur Radix UI~: sans style imposé mais avec une accessibilité garantie (navigation clavier, attributs ARIA, compatibilité lecteurs d'écran). Les composants intégrés comprennent Dialog (modales), Tabs (onglets de résultats), Progress (barre de progression d'upload), Select, Dropdown, Badge, Card, Button et Skeleton (placeholders de chargement).

\subsection{Gestion des données~: TanStack React Query 5}

\textbf{TanStack React Query} gère l'ensemble de la récupération, mise en cache et synchronisation des données provenant des APIs. Son fonctionnement est fondamental pour la fluidité de l'expérience utilisateur~:

\begin{itemize}[leftmargin=1.5em]
  \item Les données (missions, utilisateurs, analyses) sont mises en cache \textbf{2 minutes}~: naviguer entre pages ne déclenche pas d'appels réseau répétés.
  \item Après une modification (création de mission, upload de fichier), le cache est \textbf{invalidé et rechargé automatiquement en arrière-plan}, sans que l'utilisateur ait à rafraîchir la page.
  \item Les mutations gèrent automatiquement les états de chargement et d'erreur, déclenchant des notifications via la bibliothèque Sonner.
\end{itemize}

\subsection{Bibliothèques complémentaires}

Le tableau \ref{tab:frontend-libs} récapitule les bibliothèques complémentaires intégrées au projet.

\begin{table}[H]
\centering
\caption{Bibliothèques frontend complémentaires}
\label{tab:frontend-libs}
\begin{tabular}{>{\bfseries}p{4.5cm} p{9.5cm}}
\toprule
Bibliothèque & Usage \\
\midrule
Recharts 2.13 & Graphiques interactifs (camembert, histogramme, barres) basés sur D3.js avec une API React native \\
Lucide React 0.468 & Bibliothèque d'icônes SVG (tree-shaken, seules les icônes utilisées sont embarquées dans le bundle) \\
Sonner 1.7 & Notifications toast positionnées en haut à droite (succès, erreur, information, 4 secondes) \\
React Hook Form 7.54 & Gestion d'état des formulaires sans re-rendu inutile \\
Zod 3.23 & Validation de schéma à l'exécution (email \texttt{@pwc.com}, dates, formats) \\
date-fns 4.1 & Formatage des dates en français (\texttt{dd/MM/yyyy HH:mm}) \\
Axios 1.7 & Client HTTP avec intercepteurs pour l'injection automatique du token JWT \\
Prisma ORM 5.22 & Client base de données avec typage auto-généré pour PostgreSQL \\
bcryptjs 3.0 & Hachage des mots de passe (12 rounds, résistant aux attaques par force brute) \\
jose 5.9 & Génération et vérification des tokens JWT (algorithme HS256) \\
uuid 11.0 & Génération d'identifiants uniques côté client \\
\bottomrule
\end{tabular}
\end{table}

% =====================================================
\section{Architecture du Code}
% =====================================================

L'organisation du code suit les conventions Next.js App Router, structurées en couches logiques clairement séparées~:

\begin{lstlisting}[caption={Structure des dossiers du frontend}, label={lst:frontend-structure}]
frontend/src/
  app/                        # Routes et pages (App Router)
    (auth)/login/             # Page de connexion publique
    (dashboard)/              # Toutes les pages authentifiees
      dashboard/              # Tableau de bord principal
      missions/               # Liste des missions
        [id]/                 # Detail d'une mission
          analysis/           # Assistant d'analyse (wizard 4 etapes)
      history/                # Historique des analyses
      reports/                # Rapports PDF/DOCX generes
      audit-trail/            # Piste d'audit complete
      admin/users/            # Gestion des utilisateurs
    api/                      # Routes API internes Next.js
      auth/                   # Login, logout, me
      missions/               # CRUD missions
      users/                  # Gestion utilisateurs
      analysis-runs/          # Persistance des analyses
      audit-logs/             # Journal d'audit
  components/                 # Composants reutilisables
    analysis/                 # Wizard, tableau de resultats, KPI
    charts/                   # Graphiques Recharts
    datasets/                 # Upload et gestion de fichiers
    missions/                 # Cartes et formulaires de mission
    reports/                  # Generation PDF/DOCX
    explanations/             # Affichage SHAP, LIME, LLM
    layout/                   # Navbar, TopBar
    ui/                       # Primitives shadcn/ui
  lib/
    api/                      # Services d'appel aux APIs
    hooks/                    # Hooks personnalises
    store/                    # Stores locaux (mode demo)
  providers/                  # Providers React (Query, Theme, Language)
  types/                      # Interfaces TypeScript partages
  messages/                   # Traductions (fr.json, en.json)
\end{lstlisting}

% =====================================================
\section{Routes et Interfaces Utilisateur}
% =====================================================

\subsection{Page de Connexion (\texttt{/login})}

La page d'entrée est volontairement sobre et professionnelle. Elle présente le logo PwC et un formulaire à deux champs (email et mot de passe). La validation s'effectue à deux niveaux~:

\begin{itemize}[leftmargin=1.5em]
  \item \textbf{Côté client (Zod)}~: le format email est vérifié et le domaine \texttt{@pwc.com} est obligatoire. Un message d'erreur s'affiche immédiatement sans solliciter le serveur.
  \item \textbf{Côté serveur}~: le mot de passe est comparé au hash bcrypt stocké en base de données (12 rounds, environ 100\,ms par vérification --- délibérément lent pour résister aux attaques par force brute).
\end{itemize}

\subsection{Tableau de Bord Principal (\texttt{/dashboard})}

Le tableau de bord est la première page visible après la connexion. Il offre une vue synthétique de l'activité en cours~: cartes KPI (missions actives, analyses lancées, anomalies détectées), activité récente, et accès rapide aux missions. Le contenu s'adapte automatiquement au \textbf{rôle de l'utilisateur connecté}~: un auditeur ne voit que ses missions assignées, tandis qu'un manager dispose d'une vue d'ensemble complète.

\subsection{Gestion des Missions (\texttt{/missions})}

La liste des missions est présentée sous forme de cartes, chacune affichant le nom de la mission, l'entreprise cliente, le type (commissariat aux comptes, audit interne, due diligence, etc.), les dates, le statut (en préparation, en cours, terminée, archivée) et les auditeurs assignés. Une barre de recherche et des filtres combinés par statut et type permettent de naviguer rapidement.

Le bouton \textit{Nouvelle mission} --- visible uniquement pour les managers, partenaires et administrateurs --- ouvre une modale de création avec un formulaire complet incluant l'affectation d'auditeurs depuis la liste des utilisateurs actifs.

\subsection{Détail d'une Mission (\texttt{/missions/[id]})}

La page de détail consolide~:

\begin{itemize}[leftmargin=1.5em]
  \item Les \textbf{métadonnées de la mission}~: informations générales, équipe, calendrier.
  \item La \textbf{section Datasets}~: liste des fichiers CSV uploadés avec drag-and-drop (\texttt{UploadDropzone}), barre de progression, possibilité de remplacer ou supprimer un fichier (jusqu'à 500\,Mo).
  \item Le \textbf{bouton Lancer l'analyse}~: redirection vers l'assistant multi-étapes.
\end{itemize}

\subsection{Assistant d'Analyse (\texttt{/missions/[id]/analysis})}

C'est l'interface la plus sophistiquée de la plateforme. Elle guide l'utilisateur à travers quatre étapes structurées~:

\begin{description}[leftmargin=1.5em, style=nextline]
  \item[Étape 1 --- Upload du fichier CSV]
    L'utilisateur dépose ou sélectionne le fichier de transactions. Un indicateur de santé du backend ML s'affiche (modèles chargés, LLM disponible) pour rassurer l'utilisateur avant de démarrer.

  \item[Étape 2 --- Sélection du modèle]
    Trois configurations sont proposées avec leur description~:
    \begin{itemize}[leftmargin=1em]
      \item \textbf{Combiné (recommandé)}~: les trois modèles votent en ensemble pour une décision plus robuste.
      \item \textbf{XGBoost uniquement}~: rapide et précis sur les patterns de fraude connus.
      \item \textbf{AutoEncoder uniquement}~: détecte des anomalies inédites, sans supervision.
    \end{itemize}

  \item[Étape 3 --- Exécution de l'analyse]
    Une carte récapitulative s'affiche avant le lancement. Pendant l'analyse, une barre de progression et des messages d'état informent l'utilisateur de chaque phase~: détection du schéma, ingénierie des features, inférence. Le timeout est fixé à 120\,secondes.

  \item[Étape 4 --- Tableau de bord des résultats]
    Les résultats sont organisés en quatre onglets~:
    \begin{itemize}[leftmargin=1em]
      \item \textbf{Graphiques}~: camembert normal/anomalie, histogramme des scores avec ligne de seuil, barres par catégorie.
      \item \textbf{Transactions}~: tableau interactif et filtrable de toutes les transactions avec score de risque et niveau (CRITIQUE/ÉLEVÉ/FAIBLE).
      \item \textbf{Explication}~: pour une transaction sélectionnée, résumé LLM en français, valeurs SHAP, règles LIME, recommandations d'actions.
      \item \textbf{Rapport}~: génération et téléchargement du rapport en PDF ou DOCX.
    \end{itemize}
\end{description}

\subsection{Autres Pages}

\begin{itemize}[leftmargin=1.5em]
  \item \textbf{Historique (\texttt{/history})}~: liste de toutes les analyses précédentes avec filtres, permettant de retrouver et revisualiser des résultats sans relancer une analyse.
  \item \textbf{Rapports (\texttt{/reports})}~: centralise tous les rapports générés, avec téléchargement et archivage.
  \item \textbf{Piste d'audit (\texttt{/audit-trail})}~: journal complet de toutes les actions (connexions, uploads, analyses, rapports) avec horodatage et identité de l'acteur. Réservé aux managers et administrateurs.
  \item \textbf{Gestion des utilisateurs (\texttt{/admin/users})}~: création, modification et désactivation de comptes, assignation de rôles. Réservé aux administrateurs.
\end{itemize}

% =====================================================
\section{Architecture API et Communication avec le Backend}
% =====================================================

\subsection{Double couche d'API}

Le frontend communique avec deux backends distincts selon une architecture en deux niveaux~:

\begin{description}[leftmargin=1.5em, style=nextline]
  \item[Couche 1 --- Next.js API Routes (PostgreSQL)]
    Un backend léger intégré à Next.js, accessible via \texttt{/api/*}, gère toute la persistance métier~: utilisateurs, missions, datasets, historique d'analyses, piste d'audit. Ces routes s'appuient sur Prisma ORM pour interagir avec une base PostgreSQL.

  \item[Couche 2 --- FastAPI ML Backend]
    Le service Python qui héberge les modèles de machine learning, accessible via \texttt{/ml/api/*}. Le frontend ne communique jamais directement avec ce service~: les requêtes transitent par un \textbf{proxy configuré dans \texttt{next.config.ts}}, qui redirige les appels \texttt{/ml/*} vers \texttt{http://localhost:8000}. Ce mécanisme élimine les problèmes CORS et masque l'URL réelle du backend ML au navigateur.
\end{description}

\subsection{Services API}

Le tableau \ref{tab:api-ml} présente les endpoints ML consommés par le frontend, et le tableau \ref{tab:api-app} les endpoints internes.

\begin{table}[H]
\centering
\caption{Endpoints FastAPI ML consommés par le frontend}
\label{tab:api-ml}
\begin{tabular}{p{3.5cm} p{3cm} p{7cm}}
\toprule
\textbf{Service} & \textbf{Méthode / Route} & \textbf{Usage} \\
\midrule
\texttt{predict} & POST \texttt{/api/predict} & Upload CSV, pipeline complet (détection schéma, feature engineering, inférence) \\
\texttt{explain} & GET \texttt{/api/explain/\{tx\_id\}} & SHAP + LIME + explication LLM pour une transaction \\
\texttt{explainBatch} & POST \texttt{/api/explain/batch} & Explications en lot (LIME désactivé pour performance) \\
\texttt{getModels} & GET \texttt{/api/models} & Métriques des modèles (recall, précision, F1, ROC-AUC) \\
\texttt{profile} & POST \texttt{/api/profile} & Profiling qualité du dataset sans prédiction \\
\texttt{health} & GET \texttt{/api/health} & Santé du backend (modèles chargés, LLM disponible) \\
\texttt{generatePDF} & POST \texttt{/api/report} & Génération du rapport PDF \\
\texttt{generateDOCX} & POST \texttt{/api/report/docx} & Génération du rapport Word \\
\bottomrule
\end{tabular}
\end{table}

\begin{table}[H]
\centering
\caption{Endpoints Next.js internes (PostgreSQL)}
\label{tab:api-app}
\begin{tabular}{p{4cm} p{3.5cm} p{6cm}}
\toprule
\textbf{Service} & \textbf{Route} & \textbf{Usage} \\
\midrule
Auth & \texttt{/api/auth/login} & Émission du JWT \\
& \texttt{/api/auth/logout} & Invalidation de session \\
& \texttt{/api/auth/me} & Utilisateur courant \\
Missions & \texttt{/api/missions} & Liste / création \\
& \texttt{/api/missions/[id]} & Lecture / modification / suppression \\
Datasets & \texttt{/api/missions/[id]/datasets} & Liste / upload / remplacement / suppression \\
Utilisateurs & \texttt{/api/users} & Liste et détail \\
Analyses & \texttt{/api/analysis-runs} & Persistance de l'historique \\
Audit & \texttt{/api/audit-logs} & Enregistrement des actions \\
\bottomrule
\end{tabular}
\end{table}

\subsection{Intercepteurs Axios}

Les deux clients HTTP sont configurés avec des intercepteurs~:

\begin{itemize}[leftmargin=1.5em]
  \item \textbf{Intercepteur de requête}~: injecte automatiquement l'en-tête \texttt{Authorization: Bearer \{token\}} à chaque appel, en lisant le token depuis le \texttt{localStorage}.
  \item \textbf{Intercepteur de réponse}~:
    \begin{itemize}[leftmargin=1em]
      \item Réponse \textbf{401} (token expiré ou invalide) → nettoyage du \texttt{localStorage} et redirection vers \texttt{/login}.
      \item \textbf{Erreur réseau} → une tentative de relance automatique.
      \item \textbf{Erreur 4xx} (erreur client) → pas de relance.
    \end{itemize}
\end{itemize}

% =====================================================
\section{Authentification et Contrôle d'Accès}
% =====================================================

\subsection{Flux JWT}

L'authentification repose sur des \textbf{tokens JWT} (JSON Web Tokens) signés avec l'algorithme HS256. Le flux est le suivant~:

\begin{enumerate}[leftmargin=1.5em]
  \item L'utilisateur soumet son email (\texttt{@pwc.com} obligatoire) et son mot de passe.
  \item Le serveur vérifie l'email et compare le mot de passe au hash bcrypt (12 rounds).
  \item Un JWT est généré avec les informations de l'utilisateur (id, email, nom, rôle) et une durée de validité de \textbf{8 heures}.
  \item Le token est retourné au client et stocké dans le \texttt{localStorage} (\texttt{pwc\_token}).
  \item À chaque requête ultérieure, l'intercepteur Axios injecte le token dans l'en-tête \texttt{Authorization}.
  \item Si le token expire, le serveur répond 401 → déconnexion automatique et redirection vers \texttt{/login}.
\end{enumerate}

\subsection{Rôles et Permissions}

Quatre rôles définissent les droits d'accès dans la plateforme~:

\begin{description}[leftmargin=1.5em, style=nextline]
  \item[Auditeur] Accède uniquement aux missions qui lui sont assignées. Peut uploader des fichiers CSV, lancer des analyses et commenter des anomalies.
  \item[Manager] Peut créer et modifier des missions, assigner des auditeurs, approuver des rapports. Accès à toutes les missions.
  \item[Partenaire] Peut créer des missions et consulter les rapports générés.
  \item[Administrateur] Accès complet à toutes les fonctionnalités, y compris la création et la modification des comptes utilisateurs.
\end{description}

Le hook personnalisé \texttt{usePermissions()} est utilisé dans chaque composant pour afficher ou masquer conditionnellement les éléments d'interface selon le rôle de l'utilisateur connecté.

% =====================================================
\section{Déploiement des Modèles côté Frontend}
% =====================================================

\subsection{Sélection des modèles}

L'étape 2 de l'assistant d'analyse (\texttt{AnalysisWizard.tsx}) présente trois configurations de modèles. Le choix de l'utilisateur est transmis en paramètre (\texttt{model\_type}) dans le corps de la requête à \texttt{POST /ml/api/predict}, avec les valeurs possibles \texttt{combined}, \texttt{xgboost} ou \texttt{autoencoder}.

\subsection{Orchestration de l'analyse}

Quand l'utilisateur déclenche une analyse~:

\begin{enumerate}[leftmargin=1.5em]
  \item Le fichier CSV est envoyé en \textit{multipart/form-data} avec le type de modèle choisi.
  \item Une barre de progression s'affiche --- le timeout est fixé à 120 secondes pour accommoder les grands fichiers.
  \item Le backend FastAPI exécute séquentiellement~: détection du schéma, profiling qualité, ingénierie des 14 features dérivées, puis inférence des modèles sélectionnés.
  \item La réponse contient l'ensemble des transactions avec leurs scores et niveaux de risque.
  \item Si la détection de schéma est incomplète (confiance sous le seuil), un \textbf{avertissement} est affiché pour informer l'auditeur de l'impact possible sur la précision.
\end{enumerate}

\subsection{Explications à la demande}

Les explications (SHAP, LIME, résumé LLM) ne sont pas calculées en masse au moment de l'analyse~: elles sont générées \textbf{à la demande}, pour la transaction que l'auditeur sélectionne dans le tableau. Cela évite une surcharge inutile du backend pour des transactions qui ne seront jamais inspectées en détail. Pour la génération de rapport, un appel batch (\texttt{POST /ml/api/explain/batch}) permet de récupérer les explications des anomalies principales en une seule requête, avec LIME désactivé pour des raisons de performance.

% =====================================================
\section{Internationalisation et Thème}
% =====================================================

\subsection{Interface bilingue français/anglais}

L'interface est entièrement bilingue. Les traductions sont stockées dans \texttt{src/messages/fr.json} et \texttt{en.json}. La langue est persistée dans le \texttt{localStorage} (\texttt{pwc\_locale}) et peut être changée à tout moment via un sélecteur dans la barre de navigation. Le \texttt{LanguageContext} injecte la fonction de traduction \texttt{t()} dans tous les composants qui en ont besoin, avec support de l'interpolation de variables.

\subsection{Mode sombre}

Un bouton dans la barre de navigation permet de basculer entre mode clair et mode sombre. Le \texttt{ThemeProvider} gère cela via la classe \texttt{dark} sur l'élément racine HTML. TailwindCSS applique automatiquement les variantes \texttt{dark:} définies sur chaque composant.

% =====================================================
\section{Performance et Configuration}
% =====================================================

\subsection{Optimisations de build (\texttt{next.config.ts})}

\begin{itemize}[leftmargin=1.5em]
  \item \textbf{Tree-shaking des imports baril}~: seules les icônes et composants réellement utilisés de \texttt{lucide-react}, \texttt{recharts} et \texttt{@radix-ui} sont inclus dans le bundle final.
  \item \textbf{Cache agressif des assets statiques}~: 31\,536\,000 secondes (1 an) pour les images et les fichiers Next.js compilés.
  \item \textbf{Optimisation des images}~: conversion automatique en AVIF et WebP (40 à 50\,\% plus légers que JPEG/PNG).
  \item \textbf{Compression}~: activée sur toutes les réponses HTTP.
  \item \textbf{Standalone output}~: pour Docker, le build n'embarque que les dépendances strictement nécessaires à l'exécution, produisant une image légère.
  \item \textbf{Limite de 500\,Mo}~: configurée pour les Server Actions afin de permettre l'upload de fichiers CSV volumineux.
\end{itemize}

\subsection{Schéma de base de données}

Le schéma Prisma (\texttt{prisma/schema.prisma}) définit 11 modèles interconnectés~: \texttt{User}, \texttt{Mission}, \texttt{MissionAssignment}, \texttt{Dataset}, \texttt{DatasetVersion}, \texttt{AnalysisRun}, \texttt{Anomaly}, \texttt{AnomalyComment}, \texttt{Report}, \texttt{AuditLog} et \texttt{Permission}. Des index sont définis sur les colonnes les plus fréquemment filtrées (\texttt{userId}, \texttt{missionId}, \texttt{action}, \texttt{timestamp}) pour maintenir les performances de requête à mesure que les données s'accumulent.

% =====================================================
\section{Gestion des Erreurs et Expérience Utilisateur}
% =====================================================

Chaque opération susceptible d'échouer est traitée explicitement~:

\begin{itemize}[leftmargin=1.5em]
  \item \textbf{Upload échoué}~: message descriptif formaté depuis l'erreur serveur, avec possibilité de réessayer.
  \item \textbf{Analyse échouée}~: message extrait de la réponse FastAPI, avec code d'erreur si disponible.
  \item \textbf{Token expiré}~: déconnexion propre et redirection sans page blanche ni message cryptique.
  \item \textbf{Backend indisponible}~: indicateur rouge dans l'assistant d'analyse avec message d'alerte explicite.
  \item \textbf{Schéma CSV non reconnu}~: avertissement détaillé sur les colonnes manquantes et leur impact potentiel sur la précision des modèles.
\end{itemize}

Les \textbf{notifications Sonner} (4 secondes, en haut à droite) confirment toutes les actions réussies~: mission créée, fichier uploadé, analyse terminée, rapport généré.

Les \textbf{états de chargement} sont gérés avec des squelettes (\texttt{Skeleton}) qui maintiennent la mise en page pendant les requêtes, évitant les sauts visuels qui nuisent à la perception de fiabilité de l'outil.

% =====================================================
\section*{Conclusion du chapitre}
\addcontentsline{toc}{section}{Conclusion du chapitre}
% =====================================================

Le frontend de cette plateforme est une application React de niveau professionnel, pensée pour un contexte d'entreprise exigeant. Sa cohérence repose sur une architecture claire~: Next.js 15 pour le rendu hybride et le proxy ML, TanStack Query pour la gestion de l'état serveur, shadcn/ui pour l'accessibilité, et Axios avec intercepteurs pour la sécurité des échanges.

La force de cette interface réside dans la \textbf{lisibilité du parcours utilisateur}~: du premier clic sur \textit{Connexion} jusqu'au téléchargement du rapport PDF, chaque étape est guidée et protégée par des validations, des messages d'erreur clairs et des indicateurs de progression. La complexité des modèles d'intelligence artificielle est totalement abstraite~: l'auditeur choisit un fichier, sélectionne une configuration, lance l'analyse, et reçoit des résultats exploitables en quelques dizaines de secondes, accompagnés d'explications rédigées en français naturel.
